\chapter{Chick: a small language and its program repair algorithm}~\label{chick}

This chapter will focus on the development and evaluation of a tool (named
\Chick{}\footnotemark{}) whose purpose is to help functional programmers
propagate changes, made to some of their definitions, to the rest of their code.

\footnotetext{All code for \Chick{} is publicly available at:
\url{https://github.com/Ptival/chick}}

Section~\ref{chick-background} gives background specific to this chapter.

Section~\ref{chick-design} describes the global design of the tool and its
components.

Section~\ref{chick-syntax} covers the syntax of the language over which \Chick{}
operates.

Section~\ref{chick-diffs} describes a family of data types that allow us to
perform repairs.

Section~\ref{chick-lookup} will present both usual lookup rules, and
\Chick{}-specific lookup rules, that are necessary to describe the repair
algorithm.

Section~\ref{chick-repair} describes all the components of the repair algorithm.

Section~\ref{chick-deriving} goes into more details into how some diff
operations needed by the repair algorithm can be automatically derived from
small descriptions of their reactions to certain changes in the data structure
they operate on.

Section~\ref{chick-guess}, finally, describes how we automatically infer diffs
from pairs of programs without requiring additional user input.

\section{Background}

Programs are rarely written once and for all: bugs may be found that must be
fixed, requirements may evolve in ways that require rewriting parts of a
program, and programs may be re-written in semantically-equivalent ways in order
to account for performance, resource-efficiency, or even simply stylistic
concerns.

The concept of \define{refactoring}, which was introduced as early as
in~\citet{wirfs1990surveying}, characterizes program transformations that
preserve the semantics of the program being manipulated.

% TODO FIXME

\begin{figure*}[!htp]

  \noindent%
  \begin{minipage}[t]{0.50\textwidth}
    \begin{minted}[fontsize=\footnotesize,linenos=false]{coq}
Inductive nat : Set :=
| O : nat
| S : ∀ (n : nat), nat.

Inductive list (A : Type)
: Type :=
| nil : list A
| cons : A →
    list A → list A.

Definition a_list
: list nat :=
  cons nat (S O)
    (nil nat).

Definition length :
  ∀ (T : Type),
    list T → nat :=
  λ T l, list_rect T (λ _, nat) O
    (λ _ _ lt, S lt) l.

Fixpoint map :
  ∀ (A B : Type), (A → B) →
    list A →
    list B :=
  λ _ B f l,
    match l with
    | nil _      => nil B
    | cons _ h t =>
      cons B (f h) (map A B f t)
    end.
  \end{minted}
\end{minipage}%
\begin{minipage}[t]{0.50\textwidth}
  \begin{minted}[fontsize=\footnotesize,escapeinside=@@,linenos=false]{coq}
Inductive nat : Set :=
| O : nat
| S : ∀ (n : nat), nat.

Inductive @\modified{vec}@ (A : Type)
: @\modified{nat →}@ Type :=
| @\modified{vnil}@ : vec A @\modified{O}@
| @\modified{vcons}@ : A → @\modified{∀ (n : nat),}@
    @\modified{vec}@ A @\modified{n}@ → vec A @\modified{(S n)}@.

Definition a_list
: @\repaired{vec}@ nat @\repaired{(\_ : nat)}@ :=
  @\repaired{vcons}@ nat @\repaired{(\_ : nat)}@ (S O)
    (nil nat).

Definition length :
  ∀ (T : Type),
    @\repaired{vec}@ T @\repaired{(\_ : nat)}@ → nat :=
  λ T l, @\repaired{vec\_rect}@ T (λ _, nat) O
    (λ _ _ lt, S lt) l.

Fixpoint map :
  ∀ (A B : Type), (A → B) →
    @\repaired{vec}@ A @\repaired{(\_ : nat)}@ →
    @\repaired{vec}@ B @\repaired{(\_ : nat)}@ :=
  λ _ B f l,
    match l with
    | nil _         => nil B
    | @\repaired{vcons}@ _ @\repaired{\_}@ h t =>
      @\repaired{vcons}@ B @\repaired{\_}@ (f h) (map A B f t)
    end.
  \end{minted}
\end{minipage}

\caption{ We demonstrate the propagation of repairs after an inductive
definition is modified.  On the left, the original program.  On the right, we
highlight the user-modified parts in yellow.  The computed repairs are
highlighted in green. }

\label{listtovec}

\end{figure*}

\section{Design}

We will now describe the syntax of our core language, \Chick{}, that will be the
target of our repair algorithm.  In sections~\ref{chick-design-syntax-terms}
and~\ref{chick-design-syntax-programs}, we will describe the syntax of its
terms, and programs, respectively.

\subsection{Syntax of terms}\label{chick-design-syntax-terms}

\Chick{} is designed as a core functional programming language upon which the
repair algorithm will operate.  The aim was to represent a language such as
\Gallina{}, and as such, \Chick{} is a dependently-typed lambda calculus, with
inductive constructions.  The syntax of \Chick{} is given in this grammar:

\begin{grammar}%
<term> ::= \ %trick LaTeX
\alt <term> : <term>                                           \hfill (type annotation)
\alt <term> <term>                                             \hfill (application)
\alt _                                                         \hfill (hole)
\alt $\lambda$ <binder>, <term>                                \hfill (value abstraction)
\alt match <term> with $\overrightarrow{ \synt{pattern} }$ end \hfill (pattern matching)
\alt $\Pi$(<binder> : <term>) → <term>                         \hfill (type abstraction)
\alt <universe>                                                \hfill (universes)
\alt <var>                                                     \hfill (variable)

<universe> ::= Prop \alt Set \alt Type

<binder> ::= \ % trick LaTeX
\alt <var> \hfill (named)
\alt _     \hfill (anonymous)

<pattern> ::= ($\overrightarrow{ \synt{binder} }$, <term>)
\end{grammar}

The \emph{application}, \emph{value abstraction}, \emph{type abstraction}, and
\emph{variable} constructs, are standard constructs of a dependently-typed
lambda calculus.

The \emph{type annotation} construct lets us ascribe a type to a given term.
The \emph{hole} construct can take the place of a term anywhere: it stands for a
missing value that should be filled by the programmer.  Together, these two
constructs let us have \define{typed holes}, that is, placeholder values whose
type is known.  This proves invaluable throughout the repair algorithm, as we
will see, since there are situations where the algorithm must come up with
arbitrary values, knowing only the type they should bear.  In the absence of
synthesis techniques to generate such values, the algorithm can safely insert a
typed hole, leaving to the programmer the task of figuring out the proper value.

Finally, the \emph{pattern matching} construct is found in languages with
algebraic data types, and lets us dispatch code based on the shape of some
such data type.

We use the same universe names as found in \Gallina{}, though we do not yet
support its cumulative universe hierarchy syntactically.  However, this does
\emph{not} mean that \Chick{} cannot repair \Gallina{} programs that use the
universe hierarchy, but rather, it cannot syntactically represent universe
levels.  Instead, \Chick{} is oblivious to universe levels, so it will repair
programs regardless of their universe level.  The downside is that it cannot
represent explicit universe levels, and as such, might break programs that
require them.  Universe levels can be added almost for free, since the repair
algorithm could simply carry them along and attempt no repair.

While we use the underscore symbol ($\_$) to represent both a hole term and an
anonymous binder, the two concepts can never appear in the same program
location, and therefore, there can be no confusion: in a term context, the
symbol represents a hole, while in a binding context, it represents an anonymous
binder.


\subsection{Syntax of programs}\label{chick-design-syntax-programs}

Our programs are defined as sequences of commands in the following vernacular
(to borrow Coq's parlance):

\begin{grammar}
<program> ::= $\overrightarrow{ \synt{vernac} }$

<vernac> ::= \ %trick LaTeX
\alt Definition(\{ kind : <definition-kind> , name : <var>, type : <term>, body : <term> \})
\alt Inductive(<inductive>)

<definition-kind> ::= Definition | Fixpoint
\end{grammar}

Our \define{definition kinds} follow \Coq{}'s \Vernacular{} conventions: a
\mintinline{coq}{Definition} is \emph{never} recursive, while a
\mintinline{coq}{Fixpoint} is \emph{allowed} to be recursive.  Declarations are
ordered, and may only depend on previous declarations.  This will ease the task
of our repair algorithm, since dependencies are syntactically ordered.  In a
language where declarations are allowed to depend on other declarations backward
\emph{and} forward, we would want to repair programs according to the dependency
tree rather than the syntactic order of declarations.

Inductive definitions are not mutual, and follow the following syntax:

\noindent
\begin{grammar}
<inductive> ::= Inductive(\{\\
name : \synt{var},\\
parameters : $\overrightarrow{ ( \synt{var} , \synt{term} ) }$,\\
indices : $\overrightarrow{ ( \synt{binder} , \synt{term} ) }$,\\
universe : \synt{universe},\\
constructors : $\overrightarrow{ \synt{constructor} }$\\
\})

<constructor> ::= Constructor(\{\\
name : \synt{var},\\
parameters : $\overrightarrow{ ( \synt{binder}, \synt{term} ) }$,\\
indices : $\overrightarrow{ \synt{term} }$\\
\})
\end{grammar}
\vspace{\baselineskip}


\subsection{Describing program modifications with diffed data structures}\label{chick-design-diffs}

We define a family of data types that allows us to capture changes made to
programs from the language presented in
section~\ref{chick-design-syntax-programs}.  We will refer to these descriptions
of changes as \textit{diffs}.  We will usually denote a \textit{diff type}
$\delta_{\tau}$ if it corresponds to values of the diffed type $\tau$.  For
instance, the type $\delta_{\text{universe}}$ would capture how values of type
\coqinline{universe} may have been modified).

Importantly, note that we are not talking about values and types as they appear
in the user's program, but rather, about the internal values and types of the
meta-language.  For instance, if a user changes a \coqinline{bool} value in
their program from \coqinline{true} to \coqinline{false}, we will capture this
change using $\delta_{\text{term}}$, the diff type for $\synt{term}$; whereas
values of the diff type $\delta_{\text{bool}}$ would describe changes made to
boolean flags in the meta-language.

When the user makes a change to their program, we will attempt to guess the
structure of their changes, as a value of one of the program diff type
$\delta_{\text{program}}$.  To illustrate the kind of changes we are interested
in describing, let us go back to our motivating example in
Figure~\ref{listtovec}, where we can observe the following user-provided
modifications (dashed blue outline):

\begin{enumerate}

\item they renamed \coqinline{list} into \coqinline{vec},

\item added an index of type \coqinline{nat},

\item renamed constructor \coqinline{nil} into \coqinline{vnil},

\item instantiated the index for the first constructor with \coqinline{O},

\item renamed constructor \coqinline{cons} into \coqinline{vcons},

\item added a parameter \coqinline{n} of type \coqinline{nat} to the second
  constructor,

\item updated the recursive occurrence's name,

\item updated the recursive occurrence's index,

\item and instantiated the index for the second constructor with \coqinline{(S n)}.

\end{enumerate}

Intuitively, we want to capture changes like insertion, modifications,
deletions, and permutations, at all syntactic levels of the meta-language
(within terms, within inductive declarations).  We will use the same
descriptions to capture user-provided and repair-generated modifications.

Every diff type is accompanied by a patching function, which, given an element
of the diffed type and an element of the diff type produces, when successful, a
patched element.  This operation is partial for many diff types, because they
often capture modifications that only make sense for certain constructors of the
diffed type: for instance, a diff that says the head of a list has been modified
does not apply to the empty list.  We will overload the notation
$\MathPatches{x}{\delta_x}{x'}$ to indicate that $x'$ is the (optional) value
obtained when (successfully) patching $x$ according to the diff $\delta_x$,
using the relevant patching function for that diff type.  In all of our
notations, we use a black frame and a teal highlight to indicate that a value is
an output.

\subsection{Atomic diff}

There are many data types for which we will only capture changes at an atomic
granularity: either the value is the same in the new program, or it has been
replaced with a different, unrelated value.  For instance, a binder can either
have the same name, or have been renamed.  Similarly, the only possible change
to the recursive flag of a definition is to be atomically changed to a different
value.  The parameterized $\delta_{atomic}$ diff type captures such cases for a
given type $\tau$:

\begin{grammar}
<$\delta_{atomic}\ \tau$> ::= \ %trick LaTeX
\alt $\MathSame$                        \hfill (unchanged)
% \alt $\MathReplace{\langle\tau\rangle}$ \hfill (replaced)
\end{grammar}
%
with the following semantics:

\begin{mathpar}

  {
    \inferrule*
    [right=Identity]
    {  }
    {\MathPatches{x}{\MathSame}{x}}
  }

  {
    \inferrule*
    [right=Replace]
    {  }
    {\MathPatches{x}{\MathReplace{y}}{y}}
  }

\end{mathpar}

\subsection{List diff} \label{list-diff}

For lists, we provide a rich selection of diff operations.  The aim is not to
have a canonical representation, but rather to capture closely the intent of
the user modifications:

\begin{grammar}
<$\delta_{\text{list}} \ \tau \ \delta_{\tau}$> ::= \ %trick LaTeX
\alt \synt{$\delta_{\text{atomic}}\ \tau$} \hfill (atomic modification)
\alt \begin{tabular}{p{0.3cm} >{\centering}p{0.6cm} l}$\langle\tau\rangle$
       & $\MathIns{}{}$
       & \synt{$\delta_{list}\ \tau\ \delta_{\tau}$} \\\end{tabular} \hfill
     (insert a head)
\alt \begin{tabular}{p{0.3cm} >{\centering}p{0.6cm} l}$\langle\delta_\tau\rangle$
       & $\MathMod{}{}$
       & \synt{$\delta_{list}\ \tau\ \delta_{\tau}$} \\\end{tabular}
     \hfill (modify and keep the head)
%\alt \begin{tabular}{p{0.3cm} >{\centering}p{0.6cm} l}\quad & $\MathKeep{}$
%       & \synt{$\delta_{list}\ \tau\ \delta_{\tau}$} \\\end{tabular}
%     \hfill (keep the head)
\alt \begin{tabular}{p{0.3cm} >{\centering}p{0.6cm} l}\quad
       & $\MathDrop{}$
       & \synt{$\delta_{list}\ \tau\ \delta_{\tau}$} \\\end{tabular}
     \hfill (drop the head)
\alt \begin{tabular}{p{0.3cm} >{\centering}p{0.6cm} l}\quad
       & FIXME % $\MathPermute{p}{}$
       & \synt{$\delta_{\text{list}}\ \tau\ \delta_{\tau}$} \\\end{tabular}
     \hfill (permute according to a permutation $p$)
\end{grammar}
%
with the following semantics:

\begin{mathpar}
  {
    \inferrule*
    [right=Insert]
    {\MathPatches{l}{\delta_{l}}{l'}}
    {\MathPatches{l}{\MathIns{h}{\delta_{l}}}{(h :: l')}}
  }

  {
    \inferrule*
    [right=Modify]
    {\MathPatches{h}{\delta_{h}}{h'} \quad \MathPatches{t}{\delta_{t}}{t'}}
    {\MathPatches{(h :: t)}{\MathMod{\delta_{h}}{\delta_{t}}}{(h' :: t')}}
  }
\\
  % {
  %   \inferrule*
  %   [right=Keep]
  %   {\MathPatches{t}{\delta_{t}}{t'}}
  %   {\MathPatches{(h :: t)}{\MathKeep{\delta_{t}}}{(h :: t')}}
  % }
  {
    \inferrule*
    [right=Drop]
    {\MathPatches{t}{\delta_{t}}{t'}}
    {\MathPatches{(h :: t)}{\MathDrop{\delta_{t}}}{t'}}
  }

  {
    \inferrule*
    [right=Permute]
    {\MathPatches{(h_{p(1)} \Cons \ldots \Cons h_{p(|p|)} \Cons t)}{\delta}{l}}
    {\MathPatches{(h_1 \Cons \ldots \Cons h_{|p|} \Cons t)}{\MathPermute{p}{\delta}}{l}}
  }

\end{mathpar}

Note that we defined the semantics of $\MathModPiOp$ so that it both modifies
and keeps the head: the recursive occurrence in rule~\rulename{Modify},
$\delta_t$, therefore applies to the tail $t$ and not the whole list after the
head has been repaired.  On the other hand, the recursive occurrence in
rule~\rulename{Permute}, $\delta$, targets the entire list after the permutation
is performed, not solely the tail: this is necessary because we will want to
perform modifications of elements after having shuffled them around.

\subsection{Term diff}

The diffs for terms include atomic changes, as well as insertion, modification,
deletion, and permutation of most constructors.  We illustrate a couple of
these:

\begin{grammar}
<$\delta_{term}$> ::= \ %trick LaTeX
\alt \synt{$\delta_{atomic}\ t$} \hfill (atomic modification)
\alt \synt{$\delta_{term}$} $\oIns{ \$ }$ \synt{$\delta_{term}$} \hfill
(insert application)
\alt $\oIns{\lambda}$ \synt{v}, \synt{$\delta_{term}$} \hfill (insert value
abstraction)
\alt $\oIns{\Pi}$ (\synt{v} : \synt{$\delta_{term}$}),
\synt{$\delta_{term}$} \hfill (insert type abstraction)
\alt ... \hfill (other insertions)
\alt ... \hfill (removals/modifications/permutations)
\end{grammar}

\noindent Note that we use an infix dollar sign ($\$$) as a symbol for function
application in our diffs, even though we use an infix space for function
application in our terms, which is not ideal but should prevent confusion.  For
our purpose, we biased the diffs on binary operations in the least surprising
way: deleting an application keeps its left child, i.e. removes the function
call and keeps the function (see \rulename{\RmApp}), while deleting a
\coqinline{Pi} keeps its right child, i.e. removes the value being quantified
but keeps the return type (see \rulename{\RmPi}).  For insertion, we allow
maximal flexibility by passing the entire old term to all recursive occurrences:
for instance, the diff $(\MathInsApp{\MathSame}{\MathSame})$ turns any term $t$
into the self-application $(t\ t)$ (see \rulename{\InsApp}).  However, it is
often the case that only one recursive occurrence will use the original term,
while the other ones will replace it: for instance, the diff
$(\MathInsApp{\MathSame}{\MathReplace{x}})$ turns any term $t$ into the
application $(t\ x)$.

\begin{mathpar}
  {
    \inferrule*
    [right=\RmApp]
    {\MathPatches{f}{\delta}{f'}}
    {\MathPatches{f\ x}{\MathDropApp{\delta}}{f'}}
  }

  {
    \inferrule*
    [right=\RmPi]
    {\MathPatches{\tau_{2}}{\delta}{\tau_{2}'}}
    {\MathPatches{\MathPi{\tau_{1}}{x}{\tau_{2}}}{\MathDropPi{\delta}}{\tau_{2}'}}
  }

  {
    \inferrule*
    [right=\InsApp]
    {\MathPatches{t}{\delta_1}{t_1} \quad \MathPatches{t}{\delta_2}{t_2}}
    {\MathPatches{t}{\MathInsApp{\delta_1}{\delta_2}}{t_1\ t_2}}
  }

\end{mathpar}

\subsection{Other diff types}

We also need diffs for many other internal data types.  Diff types for tuples
are derived from diff types of their constituents in a straightforward way.
Inductive data type definitions, as well as constructor definitions, behave
essentially like a tuple of all their arguments, so their diff type is derived
accordingly.


\section{Repairing programs by propagating changes} \label{chick-design-repair}

In this section, we assume that we are given an original program $p$, that is, a
sequence of vernacular commands as described in
Section~\ref{chick-design-syntax-programs}, and a diff of that program
$\delta_p$ as described in Section~\ref{chick-design-diffs}, capturing a partial
refactoring made by the user.  We assume that the original proof script
type-checked, and attempt to build a repaired diff $\delta_p'$ such that:

\begin{itemize}

\item $\delta_p'$ contains the changes from $\delta_p$,

\item $\delta_p'$ completes the refactoring started by $\delta_p$, by
propagating forward the changes from $\delta_p$ to use-sites that must be
repaired to account for those changes,

\end{itemize}
%
where propagating changes forward means:

\begin{itemize}

\item propagating renaming of constants and variables,

\item propagating changes in the number, order, and type of arguments to
functions, obtained from their definition, to their use-site,

\item propagating changes in the definition of inductive data types to use-sites
of the type, its constructors, and its eliminators (those will be explained in
more details in section~\ref{deriving}).

\end{itemize}

We give a top-down description of the repair algorithm, starting at the level of
whole programs, descending all the way to terms and data type definitions.

\subsection{Repairing programs} \label{repair-program}

The repair algorithm for programs $\repairProgOp$ is described formally in
Figure~\ref{fig:repair-program} using the following judgment:

{
  \[
    {\turnstile
      {\denv{E}{\delta_E}}
      {\repairProg{p}{\delta_p}{\delta_p'}}
    }
  \]
}

\noindent To the left of the turnstile are the \textit{global environment} and
its diff.  The top part, $E$ is the environment in which the original program
$p$ was originally defined: a list of global term definitions
$\GlobalDefinitionText$, and global inductive data type definitions
$\GlobalInductiveText$.  Each vernacular command is type-checked in such a
global typing environment, and then populates it with some additional
definitions for the subsequent commands.  The bottom part, $\delta_E$, is a diff
indicating how the global environment has been changed by the time $p$ is
reached by the repair algorithm: it accounts for whether some definitions have
been added, removed, or modified.  For instance, in our motivating example, by
the time we reach the inductive definition being modified (\coqinline{list},
becoming \coqinline{vec}), $E$ would contain the global inductive data type
definition for the type \coqinline{nat}, and $\delta_E$ would keep it intact.
After this definition has been processed, $E$ would contain, in order,
\coqinline{nat} \coqinline{list}, and $\delta_E$ would still keep the former
intact, but would register the diff turning \coqinline{list} into
\coqinline{vec}.

The repair algorithm for programs takes as input the original program $p$, and
the user-provided modification $\delta_p$.  From those, it outputs a patched
modification $\delta_p'$, that propagates the changes introduced by $\delta_E$
and $\delta_p$ to the rest of the program $p$.  The algorithm essentially folds
over the sequence of vernacular commands, propagating changes from previous
commands to subsequent ones.

The rule \rulename{\RProgMod} does the bulk of the work, dispatching the repair
of each vernacular command to the repair algorithm for vernacular commands
$\repairVernacOp$ (described in Section~\ref{repair-vernacular}).  First, the
head vernacular command $v$ is repaired (using $\repairVernacOp$), returning a
repaired diff $\delta_v'$.  Then, we want to repair the rest of the program, but
the changes made to $v$ will affect the global environment for the subsequent
commands.  For instance, in our motivating example, the vernacular command
defining the inductive type \coqinline{list} gets modified to define the
inductive type \coqinline{vec} instead: when we repair the rest of the program,
we must remember that the original program was defined in a global environment
where \coqinline{list} was defined, and that the updated program must replace it
with \coqinline{vec}.  We therefore repair the rest of the program $p$ in the
appropriate updated environment and its diff.

Note that, while the rule seems to imply that vernacular commands each add
exactly one element to the global environment, it is not the case: an inductive
data type definition adds the inductive type, all its constructors, and all its
eliminators.  For simplifying diff operations, the whole $\InductiveText$
effectively acts as a placeholder for all of those when in the global
environment.  The lookup rules, given in the Appendix~\ref{appendix}, contain
the complexity of accounting for all these global definitions arising from one
inductive data type definition.

\rulename{\RProgSameCons} simply defers to \rulename{\RProgMod}: even though it
processes unchanged commands, they might need repairs to account for changes in
upstream, accounted for in the global environment.

\begin{Rules}{fig:repair-program}{ Rules for repairing programs
($\repairProgOp$).  We use $\MathSameProg$ to indicate an identity diff for
programs, and $\MathSameVernac$ to indicate an identity diff for vernacular
commands. }

  \begin{mathpar}

    {
      \inferrule*
          [lab=R-Prog-Same-Nil]
          {  }
          {\turnstile%
            { \denv{E}{\delta_E} }
            { \repairProg{[]}{\MathSame}{\MathSame} }
          }
    }

    {
      \inferrule*
          [lab=\RProgSameCons]
          {
            {\turnstile%
              { \denv{E}{\delta_E} }
              { \repairProg{\MathCons{v}{p}}{\MathMod{\MathSameVernac}{\MathSameProg}}{\delta} }
            }
          }
          {\turnstile%
            { \denv{E}{\delta_E} }
            { \repairProg{\MathCons{v}{p}}{\MathSame}{\delta} }
          }
    }

    {
      \inferrule*
          [lab=R-Prog-Replace]
          {\turnstile%
            { \denv{E}{\delta_E} }
            { \repairProg{q}{\MathSameProg}{\delta} }
          }
          {\turnstile%
            { \denv{E}{\delta_E} }
            { \repairProg{p}{\MathReplace{q}}{\delta} }
          }
    }

    % {
    %   \inferrule*
    %       [lab=\RProgKeep]
    %       {
    %         {\turnstile
    %           { \denv{E'}{\delta_E'} }
    %           { \repairProg{\MathCons{v}{p}}{\MathMod{\MathSameVernac}{\delta_p}}{\delta} }
    %         }
    %       }
    %       {\turnstile
    %         { \denv{E}{\delta_E} }
    %         { \repairProg{\MathCons{v}{p}}{\MathKeep{\delta_p}}{\delta} }
    %       }
    % }

    {
      \inferrule*
          [lab=R-Prog-Modify]
          {
            {\turnstile%
              { \denv{E}{\delta_E} }
              { \repairVernac{v}{\delta_v}{\delta_v'} }
            }
            \and
            {\turnstile%
              { \denv{v \Cons E}{\MathMod{\delta_v'}{\delta_E}} }
              { \repairProg{p}{\delta_p}{\delta_p'} }
            }
          }
          {\turnstile%
            { \denv{E}{\delta_E} }
            { \repairProg{\MathCons{v}{p}}{\MathMod{\delta_v}{\delta_p}}{\MathMod{\delta_v'}{\delta_p'}} }
          }
    }

    {
      \inferrule*
          [lab=R-Prog-Drop]
          {
            {\turnstile
              { \denv{v \Cons E}{\MathDrop{\delta_E}} }
              { \repairProg{p}{\delta_p}{\delta_p'} }
            }
          }
          {\turnstile
            { \denv{E}{\delta_E} }
            { \repairProg{\MathCons{v}{p}}{\MathDrop{\delta_p}}{\MathDrop{\delta_p'}} }
          }
    }

    {
      \inferrule*
          [lab=R-Prog-Permute]
          {
            { \MathPatches{q}{\MathPermute{p}{}}{q'} }
            \and
            {\turnstile
              { \denv{E}{\delta_E} }
              { \repairProg{q'}{\delta}{\delta'} }
            }
          }
          {\turnstile
            { \denv{E}{\delta_E} }
            { \repairProg{q}{\MathPermute{p}{\delta}}{\MathPermute{p}{\delta'}} }
          }
    }

    \end{mathpar}

\end{Rules}


\subsection{Repairing vernacular commands} \label{repair-vernacular}

We formally describe the repair algorithm for vernacular commands
$\repairVernacOp$ in Figure~\ref{fig:repair-vernac} using the following
judgment:

{
  \[
    {\turnstile
      {\denv{E}{\delta_E}}
      {\repairVernac{v}{\delta_v}{\delta_v'}}
    }
  \]
}

\noindent The rule \rulename{R-Vernac-Def} shows how we attempt to repair a
vernacular term definition.  It is quite involved because we account for the
possibility of a non-recursive definition becoming recursive and vice-versa.  We
want to account for the user-provided modifications $\delta_k$, $\delta_n$,
$\delta_\tau$, and $\delta_t$, but we might need to repair some of them.  We
first repair the type $\tau$, and obtain its repaired diff $\delta_\tau'$.  Now,
the user has modified the name of the definition according to $\delta_n$, but
there are two situations where we must intervene:

\begin{enumerate}

\item If $\delta_n$ is $\MathSame$, the user intends to keep the name $n$ for
this definition.  However, they could have introduced in the global environment,
which would be accounted for in $\delta_E$. In this case, we need to come up
with a fresh name that is free in the new environment.

\item If $\delta_n$ is $\MathReplace{m}$, the user intends to rename this
definition from $n$ into $m$.  Again, this is a problem if the new global
environment contains an $m$ already.  In this case, we also need to come up with
a fresh name that is free in the new environment.

\end{enumerate}

\noindent $\freshOneOp$ takes care of these two situations: it takes as input a
name $n$, and a desired modification for it $\delta_n$, and returns a suitable
modification $\delta_n'$, that tries to preserve $\delta_n$ when it is possible,
to preserve $n$ otherwise when it is possible, or comes up with a fresh name
when neither of these conditions can be met.

Finally, in order to repair the body of the definition, we need to build an
initial local context (section~\ref{repair-term} will discuss the local context
further).  This context is empty for non-recursive definitions, but should
contain a self-reference for recursive definitions.  The side conditions
defining $\Gamma$ and $\delta_\Gamma$ make sure that the local context is
appropriately set in all possible combinations of the recursive flag before and
after the user-modification.

We omit most of the details of repairing an inductive data type definition
behind the symbol $\repairIndOp$.  It essentially needs to go recursively in all
parameters and indices of both the inductive type itself and all of its
constructors and repair them.  Each subsequent parameter must be repaired in a
local context where its predecessor parameter has been accounted, so as to react
to possible modifications: for instance, if the first parameter must be renamed,
and the second parameter's type refers to the first parameter's name, the type
will need to be repaired.  Indices also must repaired in the context of the
repaired parameters.  Constructors are repaired one by one, in isolation, but
the repairing of their parameters.  Constructor indices are simply
instantiations of the inductive indices with terms: they are simply repaired in
isolation.

\begin{Rules}
  {fig:repair-vernac}
  { Rules for repairing vernacular commands ($\repairVernacOp$). }

  \begin{mathpar}
    {
      \inferrule*
      [lab=R-Vernac-Def]
      {
        {\MathPatches{\tau}{\delta_\tau}{\tau'}}
        \and
        {\turnstile
          { \dcontext{E}{\delta_E}{[]}{\MathSame} }
          { \repairTerm{\tau'}{\delta_\tau'}{\MathType}{\MathSame} }
        }
        \and
        {\MathPatches{\tau'}{\delta_\tau'}{\tau''}}
        \\\\\\
        {\turnstile
          {\denv{E}{\delta_E}}
          {\freshOne{n}{\delta_n}{\delta_{n}'}}
        }
        \and
        {\MathPatches{n}{\delta_{n}'}{n'}}
        \and
        {
          \Gamma \op{=}
          \begin{cases}
            []                         & \text{if } k = \text{Definition}\\
            [\MathLocalAssum{n}{\tau}] & \text{if } k = \text{Fixpoint}
          \end{cases}
        }
        \\\\\\
        {
          \delta_\Gamma \op{=}
          \begin{cases}
            \MathIns{\MathLocalAssum{n'}{\tau''}}{\MathSame} & \text{when } k = \text{Definition}\\
                                                             & \phantom{\text{when }} \delta_k = \MathReplace{\text{Fixpoint}}\\
            \MathDrop{\MathSame}                             & \text{when } k = \text{Fixpoint}\\
                                                             & \phantom{\text{when }} \delta_k = \MathReplace{\text{Definition}}\\
            \MathSame                                        & \text{otherwise}
          \end{cases}
        }
        \\\\\\
        {\MathPatches{t}{\delta_t[\subst{\delta_n}{\delta_{n}'}]}{t'}}
        \and
        {\turnstile
          {\dcontext{E}{\delta_E}{\Gamma}{\delta_\Gamma}}
          {\repairTerm{t'}{\delta_{t}'}{\tau}{\delta_\tau'}}
        }
      }
      {\turnstile
        { \denv{E}{\delta_E} }
        { \repairVernac
          {\Definition{k}{n}{\tau}{t}}
          {\ModifyDefinition{\delta_k}{\delta_n}{\delta_\tau}{\delta_t}}
          {\ModifyDefinition{\delta_k}{\delta_n'}{\delta_\tau'}{\delta_t'}}
        }
      }
    }

    {
      \inferrule*
      [lab=R-Vernac-Ind]
      {
        {\turnstile
          {\denv{E}{\delta_E}}
          {\freshOne{n}{\delta_n}{\delta_{n}'}}
        }
        \and
        {\turnstile
          {\denv{E}{\delta_E}}
          {\repairInd
            {\overrightarrow{p}}
            {\delta_{\overrightarrow{p}}}
            {\overrightarrow{i}}
            {\delta_{\overrightarrow{i}}}
            {\overrightarrow{c}}
            {\delta_{\overrightarrow{c}}}
            {(\delta_{\overrightarrow{p}}', \delta_{\overrightarrow{i}}',
              \delta_{\overrightarrow{c}}')}
          }
        }
      }
      {\turnstile
        { \denv{E}{\delta_E} }
        { \repairVernac
          {\Inductive
            {n}
            {\overrightarrow{p}}
            {\overrightarrow{i}}
            {u}
            {\overrightarrow{c}}
          }
          {\ModifyInductive
            {\delta_n}
            {\delta_{\overrightarrow{p}}}
            {\delta_{\overrightarrow{i}}}
            {\delta_u}
            {\delta_{\overrightarrow{c}}}
          }
          {\ModifyInductive
            {\delta_n'}
            {\delta_{\overrightarrow{p}}'}
            {\delta_{\overrightarrow{i}}'}
            {\delta_u}
            {\delta_{\overrightarrow{c}}'}
          }
        }
      }
    }
  \end{mathpar}

\end{Rules}


\subsection{Repairing terms} \label{chick-design-repair-term}

The repair algorithm for terms is described formally in
Figures~\ref{fig:repair-term} and~\ref{rules-unknown-repair} using the following
judgments:

{
  \[
    {\turnstile
      { \dcontext{E}{\delta_E}{\Gamma}{\delta_\Gamma} }
      {\repair{t}{\delta_t}{\tau}{\delta_\tau}}
    }
  \]
  % \[
  %   {\turnstile
  %     { \dcontext{E}{\delta_E}{\Gamma}{\delta_\Gamma} }
  %     {\genericrepair{t}{\delta_t}{\tau}}
  %   }
  % \]
  \[
    {\turnstile
      { \dcontext{E}{\delta_E}{\Gamma}{\delta_\Gamma} }
      {\repairTermWithoutType{t}{\delta_t}}
    }
  \]
}%
%
The first judgment (Figure~\ref{fig:repair-term}) describes how a term is
repaired when we know both its old type $\tau$, and how that type got modified
$\delta_\tau$.  It is syntax-directed by the structure of $\delta_\tau$.  The
second judgment (Figure~\ref{rules-unknown-repair}) describes how a term is
repaired when we no longer make decisions based on its old type and type diff:
this happens for constructs whose type does not inform us about the type of its
constituents (for instance, function application), and for simple terms (holes,
universes, variables).

The judgments for repairing terms include a local context $\Gamma$, behaving
similarly to the global environment $E$ but holding local variables introduced
by the binding structures of terms.  The first judgment is therefore to be
understood as: given that a term $t$ had type $\tau$ in the global environment
$E$ and the local context $\Gamma$, if those latter three respectively undergo
the modifications $\delta_\tau$, $\delta_E$, and $\delta_\Gamma$, then the term
$t$ should be repaired according to $\delta_t$.  The other judgment is to be
read similarly, except the type information is missing.  Importantly, the first
judgment assumes the diff $\delta_\tau$ is repaired already: it is the
responsibility of the caller to ensure so.

\input{chick-design-repair-term-diff-directed}

\input{chick-design-repair-term-term-directed}



\documentclass[preview]{standalone}

\usepackage{fontspec}
\usepackage{gfsdidot}          % makes math stuff look neat, but also overrides text stuff
\setmainfont{TeX Gyre Pagella} % forces back text stuff to look neat too
\setmonofont{DejaVuSansMono}   % monospace stuff should look neat too

\usepackage[dvipsnames]{xcolor} % must be loaded prior to tikz, pgfplots, etc.
\let\iint\relax
\let\iiint\relax
\let\iiiint\relax
\let\idotsint\relax
\usepackage{amsmath} % loaded by mathtools?
\usepackage{amssymb}
\usepackage[english]{babel}
\usepackage{booktabs}
\usepackage{calc}
\usepackage{caption}
\usepackage{color}
\usepackage{DejaVuSansMono}
\usepackage{dsfont}
\let\textlozenge\relax
\usepackage[inline]{enumitem}
\usepackage{environ}
\usepackage{etoolbox}
\usepackage{fontawesome}
\usepackage{forest}
\usepackage{graphicx}
\usepackage{longtable}
\usepackage{ltablex} \keepXColumns % Made me sad for appendix tables, where was it used?
\usepackage{makecell}
\usepackage{mathpartir}
%\usepackage{mathtools} % \coloneqq
\usepackage{mdframed}
\usepackage{minted}
\usepackage{multicol}
\usepackage[parfill]{parskip}
\usepackage{pgfplots}
\usepackage{pifont} % \ding
\usepackage{soul}
\usepackage{syntax}
\usepackage{tabularx}
\usepackage[skins,theorems]{tcolorbox}
\usepackage{tikz}
\usepackage{tikzpeople}
\usepackage{tikz-qtree}
\usepackage{xpatch}
\usepackage[breaklinks=true,pdfborder={0 0 0}]{hyperref}

% gfsdidot makes ∀ and ∃ super ugly, this reverts them to something fine
\DeclareSymbolFont{CMsymbols}{OMS}{cmsy}{b}{n}
\SetSymbolFont{CMsymbols}{bold}{OMS}{cmsy}{ub}{n}
\DeclareMathSymbol{\forall}{\mathord}{CMsymbols}{"38}
\DeclareMathSymbol{\exists}{\mathord}{CMsymbols}{"39}

\definecolor{monokaibg}{HTML}{272822}
\definecolor{color01}{HTML}{E6194B}
\definecolor{color02}{HTML}{3CB44B}
\definecolor{color03}{HTML}{FFE119}
\definecolor{color04}{HTML}{4363D8}
\definecolor{color05}{HTML}{F58231}
\definecolor{color06}{HTML}{911EB4}
\definecolor{color07}{HTML}{46F0F0}
\definecolor{color08}{HTML}{F032E6}
\definecolor{color09}{HTML}{BCF60C}
\definecolor{color10}{HTML}{FABEBE}
\definecolor{color11}{HTML}{008080}
\definecolor{color12}{HTML}{E6BEFF}
\definecolor{color13}{HTML}{9A6324}
\definecolor{color14}{HTML}{FFFAC8}
\definecolor{color15}{HTML}{800000}
\definecolor{color16}{HTML}{AAFFC3}
\definecolor{color17}{HTML}{808000}
\definecolor{color18}{HTML}{FFD8B1}
\definecolor{color19}{HTML}{000075}
\definecolor{color20}{HTML}{808080}
\definecolor{color21}{HTML}{FFFFFF}
\definecolor{color22}{HTML}{000000}

\def\mathunderline#1#2{\color{#1}\underline{{\color{black}#2}}\color{black}}

\forestset{
  default preamble={
    for tree={
      align=center,
      draw=black,
      edge={
        ultra thick,
      },
      edge path={
        \noexpand\path[\forestoption{edge}]
        (!u.parent anchor)
        -- +(0,-10pt)
        -| (.child anchor)
        \forestoption{edge label};
      },
      font=\bfseries,
      % inner sep=10pt,
      l sep=30pt,
      % line width=2pt,
      minimum height=20pt,
      minimum width=30pt,
      parent anchor=south,
      rectangle,
      ultra thick,
    }
  }
}

\newcolumntype{Y}{>{\centering\arraybackslash}X}

% defines safecoqinline, a command like coqinline but works in tabular
% environments
\makeatletter
\def\safecoqinline#1{%
\ifx\@footnotetext\TX@trial@ftn
\detokenize{#1}%
\else
\coqinline{#1}%
\fi}
\makeatother

\captionsetup{
  justification=centering,
}

\usetikzlibrary{arrows, calc, fit, positioning, shapes, tikzmark}

\tikzset{
bicolor/.style 2 args={
  dashed,dash pattern=on 4pt off 4pt,#1,
  postaction={draw,dashed,dash pattern=on 4pt off 4pt,#2,dash phase=4pt}
  },
}

\tikzstyle{Matching}=[
black,
dashed,
line width=3pt,
]

\tikzstyle{RoundedDottedPath}=[
densely dotted,
color05,
line width=3pt,
rounded corners=10pt,
]

\tikzstyle{RoundedRectangle}=[
dashed,
draw,
line width=3pt,
rounded corners=10pt,
]

\tikzstyle{NodeLabel}=[
black,
circle,
draw,
fill=white,
font=\bfseries,
inner sep=1pt,
ultra thick,
]

\graphicspath{ {./images/} }

% Fix a mdframed bug where skipbelow is ignored
% \makeatletter
% \xpatchcmd{\endmdframed}
%   {\aftergroup\endmdf@trivlist\color@endgroup}
%   {\endmdf@trivlist\color@endgroup\@doendpe}
%   {}{}
% \makeatother

% \surroundwithmdframed[
% backgroundcolor=monokaibg,
% linecolor=white,
% skipabove=1em,
% skipbelow=1em,
% leftmargin=10pt,
% innertopmargin=1pt,
% innerbottommargin=0pt
% ]{minted}

\setminted{
  % DO NOT use bgcolor, use mdframed instead
  % bgcolor=monokaibg,
  linenos=true,
  style=lovelace,
  breaklines=true,
  encoding=utf8,
  fontsize=\large,
  baselinestretch=1,
}

% need bgcolor for the display to be nice in rules
\newmintinline{coq}{bgcolor=white}

\makeatletter
\pgfdeclareshape{document}{
\inheritsavedanchors[from=rectangle] % this is nearly a rectangle
\inheritanchorborder[from=rectangle]
\inheritanchor[from=rectangle]{center}
\inheritanchor[from=rectangle]{north}
\inheritanchor[from=rectangle]{south}
\inheritanchor[from=rectangle]{west}
\inheritanchor[from=rectangle]{east}
% ... and possibly more
\backgroundpath{% this is new
% store lower right in xa/ya and upper right in xb/yb
\southwest \pgf@xa=\pgf@x \pgf@ya=\pgf@y
\northeast \pgf@xb=\pgf@x \pgf@yb=\pgf@y
% compute corner of ‘‘flipped page’’
\pgf@xc=\pgf@xb \advance\pgf@xc by-10pt % this should be a parameter
\pgf@yc=\pgf@yb \advance\pgf@yc by-10pt
% construct main path
\pgfpathmoveto{\pgfpoint{\pgf@xa}{\pgf@ya}}
\pgfpathlineto{\pgfpoint{\pgf@xa}{\pgf@yb}}
\pgfpathlineto{\pgfpoint{\pgf@xc}{\pgf@yb}}
\pgfpathlineto{\pgfpoint{\pgf@xb}{\pgf@yc}}
\pgfpathlineto{\pgfpoint{\pgf@xb}{\pgf@ya}}
\pgfpathclose
% add little corner
\pgfpathmoveto{\pgfpoint{\pgf@xc}{\pgf@yb}}
\pgfpathlineto{\pgfpoint{\pgf@xc}{\pgf@yc}}
\pgfpathlineto{\pgfpoint{\pgf@xb}{\pgf@yc}}
\pgfpathlineto{\pgfpoint{\pgf@xc}{\pgf@yc}}
}
}
\makeatother

%%%%% General-purpose text macros %%%%%
\newcommand{\define}[1]{\emph{#1}}
\newcommand{\mycite}[1]{\citeauthor{#1}~\cite{#1}}

\newcommand{\Language}[1]{\emph{#1}}
\newcommand{\Tool}[1]{\emph{#1}}
\newcommand{\Gallina}{\Language{Gallina}}
\newcommand{\Ltac}{\Language{Ltac}}
\newcommand{\Vernacular}{\Language{Vernacular}}

\newcommand{\Chick}{\Language{Chick}}
\newcommand{\Coq}{\Language{Coq}}
\newcommand{\CoqIDE}{\Tool{CoqIDE}}
\newcommand{\Haskell}{\Language{Haskell}}
\newcommand{\JavaScript}{\Language{JavaScript}}
\newcommand{\OCaml}{\Language{OCaml}}
\newcommand{\PeaCoq}{\Language{PeaCoq}}
\newcommand{\RxJS}{\Language{RxJS}}
\newcommand{\SerAPI}{\Tool{SerAPI}}
\newcommand{\Snap}{\Language{Snap}}
\newcommand{\TypeScript}{\Language{TypeScript}}

\newcommand{\OperatorColor}{purple}
\newcommand{\Operator}[1]{\textcolor{\OperatorColor}{\ #1\ }}
\newcommand{\Entails}{\Operator{\vdash}}
\newcommand{\HasType}{\Operator{:}}

%\newmintinline{coq}{fontsize=\small}
% \newcommand{\coqinline}[1]{%
%   %\colorbox{monokaibg}{%
%   \parbox[c][0.9em]{\widthof{\mycoq{#1}}}{\mycoq{#1}}%
%   %}%
% }

\newcommand{\modified}[1]{\color{Orange}{#1}}
\newcommand{\repaired}[1]{\color{PineGreen}{#1}}

\newcommand{\rulename}[1]{$\LeftTirNameStyle{#1}$}

\newcommand{\RmApp}{Rm-App}
\newcommand{\RmPi}{Rm-Pi}
\newcommand{\InsApp}{Ins-App}

%%%%% Math-mode macros %%%%%

\newcommand{\opcolor}{purple}
\newcommand{\out}[1]{ \boxed{ \textcolor{teal}{#1} } }
\newcommand{\MathPatches}[3]{%
#1 \overset{#2}{\mathbin{\textcolor{\opcolor}{\rightsquigarrow}}} \out{#3}%
}

\newcommand{\App}{\$}
\newcommand{\Mod}{Mod}
\newcommand{\Drop}{Drop}
\newcommand{\Ins}{Ins}
\newcommand{\Keep}{Keep}
\newcommand{\Lam}{\lambda}

\newcommand{\oMod}{\overset{\mathtt{\Mod}}}
\newcommand{\oDrop}{\overset{\mathtt{\Drop}}}
\newcommand{\oIns}{\overset{\mathtt{\Ins}}}
\newcommand{\oKeep}{\overset{\mathtt{\Keep}}}

\newcommand{\Cons}{::}
\newcommand{\permute}[2]{%
\overline{#2}^{\overset{#1}{\rightleftarrows}}%
}
\newcommand{\permuteOp}[2]{%
\overset{\overset{#1}{\rightleftarrows}}{#2}%
}

\newcommand{\mkMathPiRaw}[4]{#3{\Pi} #4{#1} \rightarrow #2}
\newcommand{\mkMathPi}[5]{\mkMathPiRaw{(#2 : #1)}{#3}{#4}{#5}}


\newcommand{\MathCons}[2]{#1 \Cons #2}
\newcommand{\MathDrop}[1]{\oDrop{\Cons} #1}
\newcommand{\MathDropApp}[1]{\oDrop{\App} #1}
\newcommand{\MathDropLam}[1]{\oDrop{\Lam} #1}
\newcommand{\MathDropPi}[1]{\oDrop{\Pi} #1}
\newcommand{\MathHole}{\texttt{\_}}
\newcommand{\MathIns}[2]{#1 \oIns{\Cons} #2}
\newcommand{\MathInsApp}[2]{\mkMathApp{#1}{#2}{\oIns}}
\newcommand{\MathInsAppOp}{\oIns{\App}}
\newcommand{\MathInsLam}[2]{\mkMathLam{#1}{#2}{\oIns}{}}
\newcommand{\MathInsLamOp}{\oIns{\Lam}}
\newcommand{\MathInsPi}[3]{\mkMathPi {#1}{#2}{#3}{\oIns}{}}
\newcommand{\MathInsPiOp}{\oIns{\Pi}}
\newcommand{\MathLam}[2]{\mkMathLam{#1}{#2}{}{}}
\newcommand{\MathLams}[3]{\mkMathLam{#1}{#2}{}{\BarCount{#3}}}
\newcommand{\MathKeepPiOp}{\oKeep{\Pi}}
\newcommand{\MathMod}[2]{#1 \oMod{\Cons} #2}
\newcommand{\MathModAppOp}{\oMod{\App}}
\newcommand{\MathModLamOp}{\oMod{\Lam}}
\newcommand{\MathModPiOp}{\oMod{\Pi}}
\newcommand{\MathPermute}[2]{\permuteOp{#1}{\Cons} #2}
\newcommand{\MathPermuteLams}[2]{\permuteOp{#1}{\lambda} #2}
\newcommand{\MathPermutePis}[2]{\permuteOp{#1}{\Pi} #2}
\newcommand{\MathPi}[3]{\mkMathPi{#1}{#2}{#3}{}{}}
\newcommand{\MathPis}[4]{\mkMathPi{#1}{#2}{#3}{}{\BarCount{#4}}}
\newcommand{\MathReplace}[1]{\mathds{K}({#1})}
\newcommand{\MathSame}{\mathds{1}}

\newcommand{\mkMathApp}[3]{#1 #3{ \$ } #2}

%%% REPAIR %%%

\newcommand{\op}[1]{\textcolor{\opcolor}{#1}}

\newcommand{\boundOp}{\op{\text{Bound}}}
\newcommand{\freshOneOp}{\op{\text{Fresh}_1}}
\newcommand{\freshTwoOp}{\op{\text{Fresh}_2}}
\newcommand{\repairIndOp}{\op{R_{I}}}
\newcommand{\repairProgOp}{\op{R_{P}}}
\newcommand{\repairTermOneOp}{\op{R_{T_1}}}
\newcommand{\repairTermTwoOp}{\op{R_{T_2}}}
\newcommand{\repairTermThreeOp}{\op{R_{T_2}}}
\newcommand{\repairVernacOp}{\op{R_{V}}}
\newcommand{\vernacEnvOp}{\op{E_{V}}}
\newcommand{\repairBranchesOp}{\op{R_{B}}}

\newcommand{\blackbrackets}[1]{\left[ #1 \right]}
\newcommand{\blackdiff}[2]{\blackbrackets{\genfrac{}{}{0pt}{}{#1}{#2}}}
\newcommand{\bound}[1]{\boundOp\op{(}#1\op{)}}
\newcommand{\brackets}[1]{%
  \color{\opcolor} \left[ \normalcolor #1 \color{\opcolor} \right] \normalcolor%
}
\newcommand{\context}[2]{#1 \op{,} #2}
\newcommand{\dcontext}[4]{\context{\diff{#1}{#2}}{\diff{#3}{#4}}}
\newcommand{\denv}[2]{\diff{#1}{#2}}
\newcommand{\diff}[2]{\brackets{\genfrac{}{}{0pt}{}{#1}{#2}}}
\newcommand{\dtau}[1]{\delta_{\tau_{#1}}}
\newcommand{\equals}[2]{#1 \mathbin{\textcolor{\opcolor}{=}} #2}
\newcommand{\freshOne}[3]{\freshOneOp\op{(}\diff{#1}{#2}\op{) =}\ \out{#3}}
\newcommand{\freshTwo}[3]{\freshTwoOp\op{(}#1\op{,}#2\op{) =}\ \out{#3}}
\newcommand{\genericrepair}[3]{\mkrepair{\repairTermTwoOp}{#1}{#2}{#3}{\op{?}}}
\newcommand{\nturnstile}[2]{#1 \mathbin{\textcolor{\opcolor}{\nvdash}} #2}
\newcommand{\qmark}{\textcolor{\opcolor}{?}}
\newcommand{\squiggly}[2]{#1 \mathbin{\textcolor{\opcolor}{\rightsquigarrow}} #2}
\newcommand{\subst}[2]{#1 \leftarrow #2}
\newcommand{\turnstile}[2]{#1\ \mathbin{\textcolor{\opcolor}{\vdash}}\ #2}
\newcommand{\repair}[4]{ \mkrepair{\repairTermOneOp}{#1}{#2}{#3}{#4} }
\newcommand{\repairBranches}[3]{\mksimplerepair{\repairBranchesOp}{#1\op{,}\ #2}{#3}}
\newcommand{\repairInd}[7]{
  \mksimplerepair%
  {\repairIndOp}
  {\diff{#1}{#2}\op{,}\ \diff{#3}{#4}\op{,}\ \diff{#5}{#6}}
  {#7}
}
\newcommand{\repairProg}[3]{
  \repairProgOp\op{(}\blackdiff{#1}{#2}\op{)}\ \op{=}\ \out{#3}
}
\newcommand{\repairTerm}[4]{\mkrepair{\repairTermOneOp}{#1}{#2}{#3}{#4}}
\newcommand{\repairTermWithoutType}[2]{\mksimplerepair{\repairTermThreeOp}{#1}{#2}}
\newcommand{\repairVernac}[3]{
  \repairVernacOp\op{(}\blackdiff{#1}{#2}\op{)}\ \op{=}\ \out{#3}
}

\newcommand{\DefinitionText}{\text{Definition}}
\newcommand{\InductiveText}{\text{Inductive}}
\newcommand{\ConstructorText}{\text{Constructor}}
\newcommand{\GlobalDefinitionText}{\text{GlobalDefinition}}
\newcommand{\GlobalInductiveText}{\text{GlobalInductive}}

\newcommand{\Definition}[4]{ \DefinitionText(\{ #1, #2, #3, #4 \}) }
\newcommand{\Inductive}[5]{ \InductiveText(\{ #1, #2, #3, #4, #5 \}) }
\newcommand{\Constructor}[3]{ \ConstructorText(\{ #1, #2, #3 \}) }
\newcommand{\GlobalDefinition}[3]{ \GlobalDefinitionText(\{ #1, #2, #3 \}) }
\newcommand{\GlobalInductive}[5]{ \GlobalInductiveText(\{ #1, #2, #3, #4, #5 \}) }

\newcommand{\ModifyDefinition}[4]{ \delta_{\DefinitionText}(\{ #1, #2, #3, #4 \}) }
\newcommand{\ModifyInductive}[5]{ \delta_{\InductiveText}(\{ #1, #2, #3, #4, #5 \}) }
\newcommand{\ModifyGlobalDefinition}[3]{ \delta_{\GlobalDefinitionText}(\{ #1, #2, #3 \}) }
\newcommand{\ModifyGlobalInductive}[5]{ \delta_{\GlobalInductiveText}(\{ #1, #2, #3, #4, #5 \}) }

\newcommand{\GlobalDefinitionAnon}{\GlobalDefinitionText(\{ \ldots \})}
\newcommand{\GlobalInductiveAnon}{\GlobalInductiveText(\{ \ldots \})}

\newcommand{\DefinitionAnon}{\DefinitionText(\{ \ldots \})}
\newcommand{\InductiveAnon}{\InductiveText(\{ \ldots \})}

\newcommand{\RepairProg}[1]{R-Prog-#1}

\newcommand{\RProgMod}{\RepairProg{Modify}}
\newcommand{\RProgKeep}{\RepairProg{Keep}}
\newcommand{\RProgSameCons}{\RepairProg{Same-Cons}}

\newcommand{\scalefactor}{1}
\newcommand{\invscalefactor}{1}

\NewEnviron
    {Rules}[2]
    {
      \begin{figure*}
        %\scalebox{\invscalefactor}{
          %\begin{minipage}{\scalefactor\textwidth}
            %\begin{mathpar}
              \BODY
            %\end{mathpar}
          %\end{minipage}
        %}
        \caption{#2}
        \label{#1}
      \end{figure*}
    }

\newcommand{\MathSameProg}{\MathSame_P}
\newcommand{\MathSameVernac}{\MathSame_V}
\newcommand{\MathSameGlobal}{\MathSame_G}

\newcommand{\MathProp}{\mathtt{Prop}}
\newcommand{\MathSet} {\mathtt{Set}}
\newcommand{\MathType}{\mathtt{Type}}

\newcommand{\hasType}[2]{ #1\ \op{:}\ #2 }

\newcommand{\mkrepair}[5]{
  #1
  \op{(}
  \hasType
  { #2 }
  { \diff{ #4 }{ #5 } }
  \op{) =\ }
  \out{ #3 }
}

\newcommand{\MathLocalAssum}[2]{ ( #1 : #2 ) }
\newcommand{\MathLocalDef}[3]{ ( #1 : #2 \coloneqq #3 ) }

\newcommand{\mksimplerepair}[3]{
  #1
  \op{(}
  #2
  \op{) =\ }
  \out{ #3 }
}

\newcommand{\RModPi}{\RepairTerm{Mod-$\Pi$}}
\newcommand{\RInsPi}{\RepairTerm{Ins-$\Pi$}}
\newcommand{\RDropPi}{\RepairTerm{Drop-$\Pi$}}
\newcommand{\RPermutePis}{\RepairTerm{Permute-$\Pi$s}}
\newcommand{\RReplace}{\RepairTerm{Replace}}
\newcommand{\RSame}{\RepairTerm{$\MathSame{}$-Other}}
\newcommand{\RSamePi}{\RepairTerm{$\MathSame{}$-$\Pi$}}

\newcommand{\RepairTermPrefix}{RT1}
\newcommand{\GenericRepairPrefix}{R-Term-2}
\newcommand{\UnknownTypeRepairPrefix}{R-Term-2}

\newcommand{\RepairTerm}[1]{\RepairTermPrefix-#1}
\newcommand{\GenericRepair}[1]{\GenericRepairPrefix-#1}
\newcommand{\UnknownTypeRepair}[1]{\UnknownTypeRepairPrefix-#1}

\newcommand{\mkMathLam}[4]{#3{\lambda} #4{#1} \rightarrow #2}

\newcommand{\MathModApp}[2]{ \mkMathApp{#1}{#2}{\oMod}{} }
\newcommand{\MathModLam}[2]{ \mkMathLam{#1}{#2}{\oMod}{} }
\newcommand{\MathModPi} [3]{ \mkMathPi{#1}{#2}{#3}{\oMod}{} }
\newcommand{\MathKeepPi}[1]{ \MathKeepPiOp \rightarrow #1 }

\newcommand{\MathAnnot}[2]{#1 : #2}

\newcommand{\UTRApp}{\UnknownTypeRepair{App}}
\newcommand{\UTRMatch}{\UnknownTypeRepair{Match}}
\newcommand{\UTRPi}{\UnknownTypeRepair{Pi}}
\newcommand{\UTRType}{\UnknownTypeRepair{Universe}}
\newcommand{\UTRVar}{\UnknownTypeRepair{Var}}
\newcommand{\UTRAnnot}{\UnknownTypeRepair{Annot}}
\newcommand{\UTRHole}{\UnknownTypeRepair{Hole}}
\newcommand{\UTROtherwise}{\UnknownTypeRepair{Otherwise}}

\newcommand{\repairArgsOp}{\op{R_{Fn}}}

\newcommand{\repairArgs}[5]{
  \op{\repairArgsOp(}#1\op{,} #2\op{,} #3\op{,} #4\op{) =}\ \out{#5}
}

\newcommand{\declDiff}[3]{
  \hasType
      { \diff{ #1 }{ \out{#2} } }
      { \diff{ \op{?} }{ \out{#3} } }
}

\newcommand{\MathMatch}[2]{ \text{match}\ #1\ \text{with}\ #2 }


\begin{document}

\phantom{A}
\begin{grammar}%
<term> ::= \ %trick LaTeX
\alt <var>                                                               \hfill (variable)
\alt <term> <term>                                                       \hfill (function application)
\alt \coqinline{λ} <binder> \coqinline{→} <term>                         \hfill (term abstraction)
\alt \coqinline{Π(} <binder> \coqinline{:} <term>
  \coqinline{) →} <term> \hfill (type abstraction)
\alt <universe>                                                          \hfill (universes)
\alt \coqinline{match} <term> \coqinline{with}
  $\overline{ \synt{pattern} }$ \coqinline{end} \hfill (pattern matching)
\alt <term> \coqinline{:} <term>                                         \hfill (type annotation)
\alt \coqinline{_}                                                       \hfill (hole)
%
%<universe> ::= \coqinline{Prop} | \coqinline{Set} | \coqinline{Type}
%
%<binder> ::= \ % trick LaTeX
%\alt <var> \hfill (named)
%\alt _     \hfill (anonymous)
%
%<pattern> ::= ($\overline{ \synt{binder} }$, <term>)
\end{grammar}
\phantom{A}

\end{document}

%%% Local Variables:
%%% mode: latex
%%% TeX-master: t
%%% End:

\section{Describing program modifications with diffs}\label{chick-diffs}

We define a family of data types that allows us to describe changes made to
terms from the language presented in Section~\ref{chick-syntax-terms}, as well
as programs from the language presented in Section~\ref{chick-syntax-programs}.
We will refer to these descriptions of changes as
\textit{diffs}~\footnotemark{}.  We will usually denote a \textit{diff type}
$\Delta_{\tau}$ if it corresponds to values of the diffed type $\tau$.
Informally, one can think of the type $\Delta_{\tau}$ as a descriptor for how
some value $v_{1}$ of type $\tau$ can be transformed into another value $v_{2}$.

\footnotetext{Our approach is very similar to that of~\mycite{miraldo2017type},
though it was derived independently from their work, at around the same time
they were publishing it.}

It is important to clarify how we intend to use those diff types.  We will want
to describe changes made to the data types in the abstract syntax tree
(\define{AST}) of the user's program.  An example will help us avoid some
misconception: consider the user program and its modification in
Figure~\ref{bool-modification}.  While, from the program point of view, a value
of type \coqinline{bool} was modified from the value \coqinline{True} to the
value \coqinline{False}, we will \emph{not} use the type
$\Delta_{\coqinline{bool}}$ to describe this change!  From the point of view of
the abstract syntax tree, a value of type \coqinline{term} has changed from the
original value \coqinline{Var "True"} to the value \coqinline{Var "False"},
which we will capture using the diff type $\Delta_{\coqinline{Term}}$.

To be precise, we will have to describe how a \coqinline{Definition} has
changed, and describe how its name has remained \coqinline{a}, its type has
remained \coqinline{bool}, but its definition has changed.  This will be
described as a value of type $\Delta_{\coqinline{Vernacular}}$, containing a
diff for each of those three components that may change, including the one of
diff type $\Delta_{\coqinline{Term}}$.

\begin{figure*}[!htp]

  \noindent%
  \begin{minipage}[t]{0.50\textwidth}
    \begin{minted}[escapeinside=@@,linenos=false]{coq}
Inductive bool : Set :=
| True  : bool
| False : bool.

Definition a : bool := True. @\dummystrut{ }@
\end{minted}
\end{minipage}%
\begin{minipage}[t]{0.50\textwidth}
  \begin{minted}[escapeinside=@@,linenos=false]{coq}
Inductive bool : Set :=
| True  : bool
| False : bool.

Definition a : bool := @\modified{False}@.
\end{minted}
\end{minipage}

\caption{A simple program and its modification}

\label{bool-modification}

\end{figure*}

When the user makes a change to their program, we will attempt to guess the
structure of their changes, as a value of one of the program diff type
$\Delta_{\coqinline{Program}}$.  To illustrate the kind of changes we are
interested in describing, let us look back at our motivating example in
Figure~\ref{listtovec}, where we can observe the following partial attempts at
refactoring:

\begin{enumerate}

\item they renamed \coqinline{list} into \coqinline{vec},

\item added an index of type \coqinline{nat},

\item renamed constructor \coqinline{nil} into \coqinline{vnil},

\item instantiated the index for the first constructor with \coqinline{O},

\item renamed constructor \coqinline{cons} into \coqinline{vcons},

\item added a parameter \coqinline{n} of type \coqinline{nat} to the second
  constructor,

\item updated the recursive occurrence's name,

\item updated the recursive occurrence's index,

\item and instantiated the index for the second constructor with \coqinline{(S n)}.

\end{enumerate}

Intuitively, we want to capture changes like insertion, modifications,
deletions, and permutations, at all syntactic levels of the meta-language
(within terms, within inductive declarations, etc.).  We will use the same
descriptions to capture user-provided and repair-generated modifications.

Every diff type is accompanied by a patching function.  Given an element of the
diffed type, say $\MathLocalAssum{x}{\chi}$, and an element of the diff type for
that particular type, say $\MathLocalAssum{\delta_{x}}{\Delta_{\chi}}$, the
patching function produces, when successful, a patched element (with the same
type as the original one), say $\MathLocalAssum{x'}{\chi}$.  This operation is
partial for many diff types, because they often capture modifications that only
make sense for certain constructors of the diffed type: for instance, a diff
stating that the head of a list has been modified can not be meaningfully
applied to an empty list.  We will overload the notation
$\MathPatches{x}{\delta_x}{x'}$ to indicate that $x'$ is the (optional) value
obtained when (successfully) patching $x$ according to the diff $\delta_x$,
using the relevant patching function for that diff type.  In all of our
notations, we use a black frame and a teal highlight to indicate that a value is
an output.

\subsection{Atomic diff}

There are many data types for which we will only capture changes at an atomic
granularity: either the value is the same in the new program, or it has been
replaced with a different, unrelated value.  For instance, a binder can either
have the same name, or have been renamed.  Similarly, the only possible change
to the recursive flag of a definition is to be atomically changed to a different
value.  The parameterized diff type $\DeltaAtomic{}$ captures such cases
for a given type $\tau$:

\begin{grammar}
<$\DeltaAtomic{}\ \tau$> ::= \ %trick LaTeX
\alt $\MathSame$                        \hfill (unchanged)
\alt $\MathReplace{\langle\tau\rangle}$ \hfill (replaced)
\end{grammar}
%
with the following semantics:

\begin{mathpar}
  {
    \inferrule*
    [right=Identity]
    {  }
    {\MathPatches{x}{\MathSame}{x}}
  }

  {
    \inferrule*
    [right=Replace]
    {  }
    {\MathPatches{x}{\MathReplace{y}}{y}}
  }

\end{mathpar}

\paragraph{Notation} We chose this notation to be reminiscent of the identity
function (for $\MathSame{}$), and the constant function (for
$\MathReplaceOp{}$).  However, these are not functions, but simply inert
constructors.

\subsection{List diff}~\label{list-diff}

We will now describe our diff type for lists.  \emph{Again}, the subject of
discourse is \emph{not} lists as they appear in the user's program, but lists of
abstract syntax tree constructs, as they appear in the AST of said program.  We
refer the reader to Section~\ref{chick-syntax-terms} and
Section~\ref{chick-syntax-programs}, to remind themselves of the nature of such
lists and their prevalence: the patterns of a \coqinline{match} construct, the
parameters and indices to an inductive type, as well as programs themselves, are
all instances of lists of syntactic constructs.

We provide a rich selection of diff operations for lists.  The aim is not to
have a canonical representation, but rather to capture closely the intent of the
user modifications.  We give the abstract syntax for these operators as prefix
and infix operators, following a visual intuition of what happens to a list
\coqinline{h ∷ t}, though the reader is encouraged to read the semantics,
immediately following, in order to understand the intent of each construct, as
it will probably feel opaque on first glance.  We also provide a concrete
example of a diff in Section~\ref{chick-diffs-example} that uses many of these
constructs.%
%
\footnote{Note that we omit the constructor for the atomic list diff, even
though our implementation requires it to lift values of the atomic diff type to
the corresponding list diff type.}.

\begin{grammar}
<$\DeltaList{}\ \tau \ \Delta_{\tau}$> ::= \ %trick LaTeX
\alt \synt{$\DeltaAtomic{}\ \tau$} \hfill (atomic modification of the whole list)
\alt \begin{tabular}{p{0.9cm} >{\centering}p{0.9cm} l}$\langle\tau\rangle$
       & $\MathIns{}{}$
       & \synt{$\DeltaList{}\ \tau\ \Delta_{\tau}$} \\\end{tabular} \hfill
     (insert a head)
\alt \begin{tabular}{p{0.9cm} >{\centering}p{0.9cm} l}$\langle\Delta_\tau\rangle$
       & $\MathMod{}{}$
       & \synt{$\DeltaList{}\ \tau\ \Delta_{\tau}$} \\\end{tabular}
     \hfill (modify and keep the head)
\alt \begin{tabular}{p{0.9cm} >{\centering}p{0.9cm} l}\quad
       & $\MathDrop{}$
       & \synt{$\DeltaList{}\ \tau\ \Delta_{\tau}$} \\\end{tabular}
     \hfill (drop the head)
\alt \begin{tabular}{p{0.9cm} >{\centering}p{0.9cm} l}\quad
       & $\MathPermute{p}{}$
       & \synt{$\DeltaList{}\ \tau\ \Delta_{\tau}$} \\\end{tabular}
     \hfill (permute according to a permutation $p$)
\end{grammar}
%
with the following semantics:

\begin{mathpar}
  {
    \inferrule*
    [right=Insert]
    {\MathPatches{l}{\delta_{l}}{l'}}
    {\MathPatches{l}{\MathIns{h}{\delta_{l}}}{(h :: l')}}
  }

  {
    \inferrule*
    [right=Modify]
    {\MathPatches{h}{\delta_{h}}{h'} \quad \MathPatches{t}{\delta_{t}}{t'}}
    {\MathPatches{(h :: t)}{\MathMod{\delta_{h}}{\delta_{t}}}{(h' :: t')}}
  }
\\
  % {
  %   \inferrule*
  %   [right=Keep]
  %   {\MathPatches{t}{\delta_{t}}{t'}}
  %   {\MathPatches{(h :: t)}{\MathKeep{\delta_{t}}}{(h :: t')}}
  % }
  {
    \inferrule*
    [right=Drop]
    {\MathPatches{t}{\delta_{t}}{t'}}
    {\MathPatches{(h :: t)}{\MathDrop{\delta_{t}}}{t'}}
  }

  {
    \inferrule*
    [right=Permute]
    {\MathPatches{(h_{p(1)} \Cons \ldots \Cons h_{p(|p|)} \Cons t)}{\delta}{l}}
    {\MathPatches{(h_1 \Cons \ldots \Cons h_{|p|} \Cons t)}{\MathPermute{p}{\delta}}{l}}
  }

\end{mathpar}

Note that we defined the semantics of $\MathModPiOp$ so that it both modifies
and keeps the head: the recursive occurrence in rule \rulename{Modify},
$\delta_t$, therefore applies to the tail $t$ and not the whole list after the
head has been repaired.  On the other hand, the recursive occurrence in rule
\rulename{Permute}, named $\delta$, targets the entire list after the
permutation is performed, not solely the tail: this is necessary because we will
want to perform modifications of elements after having shuffled them around.

\subsection{Term diff}\label{chick-diffs-term-diff}

The diffs for terms include atomic changes, as well as insertion, modification,
deletion, and permutation of most constructors.  We illustrate a couple of
these:

\begin{grammar}
<$\DeltaTerm{}$> ::= \ %trick LaTeX
\alt \synt{$\DeltaAtomic{}\ t$} \hfill (atomic modification)
\alt \synt{$\DeltaTerm{}$} $\oIns{ \$ }$ \synt{$\DeltaTerm{}$} \hfill
(insert application)
\alt $\oIns{\lambda}$ \synt{v}, \synt{$\DeltaTerm{}$} \hfill (insert value
abstraction)
\alt $\oIns{\Pi}$ (\synt{v} : \synt{$\DeltaTerm{}$}),
\synt{$\DeltaTerm{}$} \hfill (insert type abstraction)
\alt … \hfill (other insertions)
\alt … \hfill (removals/modifications/permutations)
\end{grammar}

\noindent Note that we use an infix dollar sign ($\$$) as a symbol for function
application in our diffs, even though we use an infix space for function
application in our terms, which is not ideal, but should help with readability.
For our purpose, we biased the diffs on binary operations in the least
surprising way: deleting a function application keeps its left child, i.e.\
removes the function call and keeps the function (Rule \rulename{\RmApp}), while
deleting a $\Pi$ keeps its right child, i.e.\ removes the value being quantified
but keeps the return type (Rule \rulename{\RmPi}).  For insertion, we allow
maximal flexibility by passing the entire old term to both recursive
occurrences: for instance, the diff $(\MathInsApp{\MathSame}{\MathSame})$ turns
any term $t$ into the self-application $(t\ t)$ (Rule \rulename{\InsApp}).
However, it is often the case that only one recursive occurrence will use the
original term, while the other ones will replace it: for instance, the diff
$(\MathInsApp{\MathSame}{\MathReplace{x}})$ turns any term $t$ into the
application $(t\ x)$.

\begin{mathpar}
  {
    \inferrule*
    [right=\RmApp]
    {\MathPatches{f}{\delta}{f'}}
    {\MathPatches{f\ x}{\MathDropApp{\delta}}{f'}}
  }

  {
    \inferrule*
    [right=\RmPi]
    {\MathPatches{\tau_{2}}{\delta}{\tau_{2}'}}
    {\MathPatches{\MathPi{\tau_{1}}{x}{\tau_{2}}}{\MathDropPi{\delta}}{\tau_{2}'}}
  }

\\

  {
    \inferrule*
    [right=\InsApp]
    {\MathPatches{t}{\delta_1}{t_1} \quad \MathPatches{t}{\delta_2}{t_2}}
    {\MathPatches{t}{\MathInsApp{\delta_1}{\delta_2}}{t_1\ t_2}}
  }

\end{mathpar}

\subsection{Other diff types}

We also need diffs for many other internal data types.  Diff types for tuples
are derived from diff types of their constituents in a straightforward way.
Inductive data type definitions, as well as constructor definitions, behave
essentially like a tuple of all their arguments, so their diff type is derived
accordingly.

\subsection{Example of a program diff}~\label{chick-diffs-example}

With all of this machinery, we can define the original diff for our running
example (again, referring to the changes seen on Figure~\ref{listtovec}).  The
diff is a value of type $\DeltaInductive{}$, as shown in
Figure~\ref{diff-list-vec}.  Because the diffs for the constructors take a lot
of space, they are abbreviated as $\delta \coqinline{nil}$ and $\delta
\coqinline{cons}$, and shown separately in Figures~\ref{diff-list-vec-nil} for
\coqinline{nil} and~\ref{diff-list-vec-cons} for \coqinline{cons}.

\begin{figure*}[!htp]

  \centering

  We demonstrate a diff between the following two programs.  The constructors
are elided, and their diff ($\delta\coqinline{nil}$ and
$\delta\coqinline{cons}$) is shown in Figures~\ref{diff-list-vec-nil}
and~\ref{diff-list-vec-cons} respectively.

\noindent%
\cprotect\fbox{%
  \begin{minipage}[t]{0.50\textwidth-2\fboxrule-2\fboxsep}%
    \begin{minted}[fontsize=\footnotesize,escapeinside=@@,linenos=false]{coq}
Inductive nat : Set :=
| O : nat
| S : ∀ (n : nat), nat.

Inductive list (A : Type) @\dummystrut{ }@
: Type := …               @\dummystrut{ }@

… (* rest of first program *)
    \end{minted}
  \end{minipage}}%
\cprotect\fbox{%
  \begin{minipage}[t]{0.50\textwidth-2\fboxrule-2\fboxsep}%
    \begin{minted}[fontsize=\footnotesize,escapeinside=@@,linenos=false]{coq}
Inductive nat : Set :=
| O : nat
| S : ∀ (n : nat), nat.

Inductive @\modified{vec}@ (A : Type)
: @\modified{nat →}@ Type := …

… (* identical *)
    \end{minted}
  \end{minipage}}

  \vspace{2em}%

  \begin{align*}
&\MathSameInductive{} \MathMod{}{} && \text{do not modify inductive \coqinline{nat}}\\
& \deltaInductive{} ( && \text{modify \coqinline{list} into \coqinline{vec}} \\
& \qquad \MathReplace{\coqinline{vec}}, && \text{modify the name} \\
& \qquad \MathSameList{},       && \text{keep the parameter} \\
& \qquad \MathIns{\coqinline{nat}}{\MathSameList{}}, && \text{add the \coqinline{nat} index} \\
& \qquad \MathSameUniverse{}, && \text{keep the universe} \\
& \qquad \MathMod{\delta\coqinline{nil}}{\MathMod{\delta\coqinline{cons}}{\MathSameList{}}}
  && \text{modify the constructors (elided)} \\
& ) \MathMod{}{} && \\
& \MathSameList{} && \text{do not modify the rest of the program}
  \end{align*}

  \caption{Diff for our running example (constructors elided)}
  \label{diff-list-vec}

\end{figure*}

While the inductive diff should be somewhat straightforward, the reader might be
surprised by $\delta\coqinline{nil}$ (in Figure~\ref{diff-list-vec-nil})
\emph{not} mentioning the renaming of \coqinline{list} (on the left) into
\coqinline{vec} on the right.  While this renaming appears syntactically in the
concrete syntax, it is only a surface-level syntactic requirement that
constructors repeat the type they inhabit.  The abstract syntax tree only
contains the list of indices, and does \emph{not} repeat the name of the inductive
type, nor the parameters of the inductive type, in every constructor.

\begin{figure*}[!htp]

  \noindent%
\cprotect\fbox{%
  \begin{minipage}[t]{0.50\textwidth-2\fboxrule-2\fboxsep}%
    \begin{minted}[fontsize=\footnotesize,escapeinside=@@,linenos=false]{coq}
…
| nil : list A @\dummystrut{ }@
  \end{minted}
\end{minipage}}%
\cprotect\fbox{%
  \begin{minipage}[t]{0.50\textwidth-2\fboxrule-2\fboxsep}%
  \begin{minted}[fontsize=\footnotesize,escapeinside=@@,linenos=false]{coq}
…
| @\modified{vnil}@ : @\modified{vec}@ A @\modified{O}@
  \end{minted}
\end{minipage}}

  \vspace{2em}%

  \begin{align*}
& \delta\coqinline{nil} \coloneqq & && \\
& \qquad \deltaConstructor{} ( && \text{modify the \coqinline{nil} constructor} \\
& \qquad \qquad \MathReplace{\coqinline{vnil}}, && \text{modify constructor name} \\
& \qquad \qquad \MathSame{},                    && \text{no change to parameters} \\
& \qquad \qquad \MathIns{\coqinline{O}}{\MathSame{}}
  && \text{instantiate the \coqinline{nat} index with value \coqinline{O}} \\
& \qquad ) && \\
  \end{align*}

  \caption{Diff for our running example (\coqinline{nil} constructor only)}
  \label{diff-list-vec-nil}

\end{figure*}

\begin{figure*}[!htp]

  \noindent%
\cprotect\fbox{%
  \begin{minipage}[t]{0.50\textwidth-2\fboxrule-2\fboxsep}%
    \begin{minted}[fontsize=\footnotesize,escapeinside=@@,linenos=false]{coq}
…
| cons : A →         @\dummystrut{ }@
    list A → list A. @\dummystrut{ }@
  \end{minted}
\end{minipage}}%
\cprotect\fbox{%
  \begin{minipage}[t]{0.50\textwidth-2\fboxrule-2\fboxsep}%
  \begin{minted}[fontsize=\footnotesize,escapeinside=@@,linenos=false]{coq}
…
| @\modified{vcons}@ : A → @\modified{∀ (n : nat),}@
    @\modified{vec}@ A @\modified{n}@ → @\modified{vec}@ A @\modified{(S n)}@.
  \end{minted}
\end{minipage}}

  \vspace{2em}%

  \begin{align*}
& \delta\coqinline{cons} \coloneqq & && \\
& \qquad \delta_{\ConstructorText} (  & && \text{modify the \coqinline{cons} constructor} \\
& \qquad \qquad \MathReplace{\coqinline{vcons}}, & && \text{modify constructor name} \\
& \qquad \qquad \MathSame{} & \MathMod{}{} && \text{keep the first parameter} \\
& \qquad \qquad ( \coqinline{n} : \coqinline{nat} ) & \MathIns{}{} && \text{insert a second parameter} \\
& \qquad \qquad ( \MathSameBinder{} : \MathInsApp
  {(\MathModApp{\MathReplace{\coqinline{vec}}}{\MathSameTerm{}})}
  {\coqinline{n}}
) & \MathMod{}{}
  && \text{modify \coqinline{list A} into \coqinline{vec A n}} \\
& & && \text{while keeping binder anonymous} \\
& \qquad \qquad \MathSameList{}, & && \text{no other parameter} \\
& \qquad \qquad \coqinline{(S n)} & \MathIns{}{} && \text{instantiate the \coqinline{nat} index} \\
& \qquad \qquad \MathSameList{}, & && \text{no other index} \\
& \qquad ) && \\
  \end{align*}
  \caption{Diff for our running example (\coqinline{cons} constructor only)}\label{diff-list-vec-cons}
\end{figure*}

One might also find the modification of the parameter in
Figure~\ref{diff-list-vec-cons} hard to read.  The diff $\MathInsApp
{(\MathModApp{\MathReplace{\coqinline{vec}}}{\MathSameTerm{}})} {\coqinline{n}}$
might seem inverted, but it only appears so because function application is
left-associative.  Therefore, when comparing \coqinline{list A} with
\coqinline{vec A n}, we are comparing \coqinline{(app list A)} and
\coqinline{(app (app vec A) n)}.  In this form, we can better see that the
outermost application of \coqinline{(app list A)} corresponds to the innermost
application of \coqinline{(app (app vec A) n)}, that is, \coqinline{(app vec
A)}.  This should make it clear that the diff between these two terms needs to
insert an application node first (the one with argument \coqinline{n}), and then
will modify the leftmost argument from \coqinline{list} to \coqinline{vec}.  The
reader might want to refer to rule \rulename{\InsApp} from
Section~\ref{chick-diffs-term-diff} for the semantics of $\MathInsApp{}{}$.

Again, the diff shown in Figure~\ref{diff-list-vec} would be the input to our
repair algorithm, alongside the original program.  The repair algorithm should
propagate changes in such a way as to generate all the other fixes seen in teal
in Figure~\ref{listtovec}.

\subsection{Why compute diffs rather than repaired
values?}~\label{chick-diffs-importance}

We have now hinted multiple times to the importance of computing and storing
diffs for our repair algorithm, rather than just computing and storing the
repaired data.  We now have the necessary syntax and semantics to explain how
the two differ.  Consider an original function \coqinline{f}, with type
\coqinline{nat → bool → string}.  Now consider the following two diffs that can
change this type to \coqinline{bool → nat → string}:

\begin{align*}
\delta_{1} & = \MathPermutePis{[1, 0]}{\MathSame{}}\\
\delta_{2} & =
  \MathDropPi{(\MathModPi
             {\MathSame{}}
             {\MathSame{}}
             {\MathInsPi{\coqinline{_}}{\coqinline{nat}}{\MathSame{}}}
             )}
\end{align*}

\noindent%
%
Indeed if we run our patching function, we obtain:

\begin{align*}
\MathPatches
{\coqinline{nat → bool → string}}
{\delta_{1}}
{\coqinline{bool → nat → string}}\\
\MathPatches
{\coqinline{nat → bool → string}}
{\delta_{2}}
{\coqinline{bool → nat → string}}
\end{align*}

Now, consider some client code that calls this function \coqinline{f}.  For
instance, prior to repairing, let's say that we have a call:

\begin{minted}[linenos=false]{coq}
let result := f (S n) b in …
\end{minted}

If the only information we have available, about the type of \coqinline{f}, is
that its new type is \coqinline{bool → nat → string}, all we can do is make a
wild guess as to how we should modify this function call:

\begin{itemize}
\item should the \coqinline{bool} argument be \coqinline{b}?
\item should the \coqinline{nat} argument be \coqinline{(S n)}?
\end{itemize}

Knowing which diff turned the old type of \coqinline{f} into the new type gives
us much more structural information, allowing us to make an educated guess.  If
the diff was $\delta_{1}$, that is, $\MathPermutePis{[1, 0]}{\MathSame{}}$, then
it appears the arguments were swapped, and so we should patch the client code as
follows:

\begin{minted}[linenos=false]{coq}
let result := f b (S n) in …
\end{minted}

\noindent%
%
On the other hand, if the diff was $\delta_{2}$, that is,
$\MathDropPi{(\MathModPi{\MathSame{}}{\MathSame{}}
{\MathInsPi{\coqinline{_}}{\coqinline{nat}}{\MathSame{}}})}$, then it appears
that the parameter of type \coqinline{bool} has moved to first position, but the
parameter of type \coqinline{nat} is a brand new parameter, unrelated to the one
that was present in the old type.  Therefore, the patch to the client code
should rather produce the following:

\begin{minted}[linenos=false]{coq}
let result := f b (_ : nat) in …
\end{minted}

\noindent where we inserted a typed hole of type \coqinline{nat} rather than
conserve the value \coqinline{(S n)} from the old program, as our diff indicates
it is not relevant.

This should illustrate how diffs contain more information than before-after
pairs.  As long as we keep propagating changes through the program as diffs,
rather than patched constructs, we can retain the semantic intent of the changes
throughout our repair algorithm.  However, this relies on obtaining an original
diff, based on the refactoring attempt from the user, that captures their
intent.  We will describe how we can guess satisfactory diffs in
Section~\ref{chick-guess}, though we will sometimes need the programmer to
disambiguate their intent explicitly.

\section{Lookup rules}~\label{chick-lookup-rules}

When encountering a variable in a program under repair, we will often need to
answer the question ``What has happened to this variable?''.  In order to do so,
we generalize the usual notion of variable lookup.  In presentations of lambda
calculi where both a global environment and a local context carry typing
information, the usual notion of variable lookup first looks in the local
context (since it carries the ``closest'' enclosing scopes), then, if no binding
is found, looks in the global environment.


We use the following judgments for looking up a variable in the local context
$\Gamma$ (on the left), and in the global environment $E$ (on the
right):~\footnote{While the two judgments share the same syntax, it should
always be clear whether the value to the left of the turnstile is a local
context or a global environment.}

\[
  \turnstile{\Gamma}{ \hasType{ v }{ \out{\tau} } }
  \hspace{1cm}
  \turnstile{E}{ \hasType{ v }{ \out{\tau} } }
\]

\subsection{Lookup rules for types}

The rules for looking up types are straightforward, as shown in
Figures~\ref{lookup-type-dispatch},~\ref{lookup-type-context},
and~\ref{lookup-type-environment}.  In the latter, we handle the complexity of
inductive data type definitions.  For those, a system like \Coq{} will introduce
not only a name for the inductive type being defined, but also one name per
elimination principles it generates, and finally one name per the constructor of
the data type.  This is captured in the lookup rules, where we check those in
order before looking in the rest of the environment.  Unfortunately,
rule~\rulename{\LookupEnv-There} has to be quite bloated, but its premises only
ensure that we look in the rest of the environment only when neither of the
three previous rules hold.

\begin{Rules}{lookup-type-dispatch}{ Lookup rules (local context and global environment) }

\begin{mathpar}
  {
    \inferrule*
    [lab=\LookupCtxt]
    { \turnstile{ \Gamma }{ \hasType{ v }{ \out{\tau} } } }
    { \turnstile{ \context{E}{\ \Gamma} }{ \hasType{ v }{ \out{\tau} } } }
  }

  {
    \inferrule*
    [lab=\LookupEnv]
    {
      \nturnstile{ \Gamma }{ \hasType{ v }{ \out{\tau} } }
      \and
      \turnstile{ E }{ \hasType{ v }{ \out{\tau} } }
    }
    { \turnstile{ \context{E}{\ \Gamma} }{ \hasType{ v }{ \out{\tau} } } }
  }

\end{mathpar}

\end{Rules}

\begin{Rules}{lookup-type-context}{ Lookup rules (local context) }

\begin{mathpar}
  {
    \inferrule*
    [lab=\LookupCtxt-Here]
    {  }
    { \turnstile{ \MathCons{\MathLocalAssum{v}{\tau}}{\Gamma} }{ \hasType{ v }{ \out{\tau} } } }
  }

  {
    \inferrule*
    [lab=\LookupCtxt-There]
    { w \neq v \and \turnstile{ \Gamma }{ \hasType{ v }{ \out{\tau} } } }
    { \turnstile{ \MathCons{\MathLocalAssum{w}{\tau}}{\Gamma} }{ \hasType{ v }{ \out{\tau} } } }
  }

\end{mathpar}

\end{Rules}

\begin{Rules}{lookup-type-environment}{ Lookup rules (global environment) }

\begin{mathpar}

  {
    \inferrule*
    [lab=\LookupEnv-Definition-Here]
    {  }
    { \turnstile{ {\MathCons{\Definition{k}{v}{\tau}{t}}{E}} }{ \hasType{ v }{ \out{\tau} } } }
  }

  {
    \inferrule*
    [lab=\LookupEnv-Definition-There]
    { w \neq v \and \turnstile{ E }{ \hasType{ v }{ \out{\tau} } } }
    { \turnstile{ {\MathCons{\Definition{k}{w}{\tau}{t}}{E}} }{ \hasType{ v }{ \out{\tau} } } }
  }

  {
    \inferrule*
    [lab=\LookupEnv-Inductive-Type-Here]
    {
      {
        I = \Inductive%
        {n_{\text{ind}}}
        {\overline{p_{\text{ind}}}}
        {\overline{i_{\text{ind}}}}
        {u}
        {\overline{c}}
      }
      \\
      {
        \text{either}
        \begin{cases}
          v = n_{ind}
          & \text{and\ } \op{\text{InductiveType}(} I \op{)} = \out{\tau} \\
          \text{or}\\
          \exists u_{\text{elim}} \text{\ such that\ } \\
          \quad v = \op{\text{EliminatorName}(} I \op{,\ } u_{\text{elim}} \op{)}
          & \text{and\ }
          \op{\text{EliminatorType}(} I \op{,\ } u_{\text{elim}} \op{)} = \out{\tau} \\
          \text{or}\\
          \exists\ C \in \overline{c} \text{\ such that\ } \\
          \quad C = \Constructor%
          {v}
          {\overline{p_{\text{ctor}}}}
          {\overline{i_{\text{ctor}}}}
          & \text{and\ }
          \op{\text{ConstructorType}(} I \op{,\ } C \op{)}
          = \out{\tau}
        \end{cases}
      }
    }
    { \turnstile{ {\MathCons{I}{E}} }{ \hasType{ v }{ \out{\tau} } } }
  }

  {
    \inferrule*
    [lab=\LookupEnv-There]
    {
      {
        I = \Inductive%
        {n_{\text{ind}}}
        {\overline{p_{\text{ind}}}}
        {\overline{i_{\text{ind}}}}
        {u}{\overline{c}}
      }
      \\
      {
        \text{all of}
        \begin{cases}
          v \neq n_{\text{ind}} & \\
          \text{and}\\
          \forall u_{\text{elim}},
          v \neq \op{\text{EliminatorName}(} I \op{,\ } u_{\text{elim}} \op{)}\\
          \text{and}\\
          \forall \Constructor%
          {n_{\text{ctor}}}
          {\overline{p_{\text{ctor}}}}
          {\overline{i_{\text{ctor}}}}
          \in \overline{c}, v \neq n_{\text{ctor}}
        \end{cases}
      }
      \\\\
      { \turnstile{ E }{ \hasType{v}{\out{\tau}} } }
    }
    { \turnstile{ {\MathCons{I}{E}} }{ \hasType{ v }{ \out{\tau} } } }
  }

\end{mathpar}

\end{Rules}

\subsection{Lookup rules for diffs}

Our second notion of a lookup essentially follows the same strategy, but we are
\emph{\textbf{not}} interested in finding the type of the variable, but rather
in knowing \textbf{how} \emph{both} the binding \emph{and} the binding type have
been modified.  To do so, we use the following judgment:

{
  \[
    \turnstile
    {\dcontext {E} {\delta_E} {\Gamma} {\delta_\Gamma} }
    { \declDiff { v } { \delta_v } { \delta_{\tau} } }
  \]
}

\paragraph{Notation} This judgment can be read as follows:
\begin{itemize}

\item \textbf{Top}: If you are interested in the variable $v$, as bound in the
original global environment $E$ and original local context $\Gamma$,

\item \textbf{Bottom left}: and the global environment underwent modification
$\delta_{E}$,

\item \textbf{Bottom left}: and the local context underwent modification
$\delta_{\Gamma}$,

\item \textbf{Bottom right}: then, references to the variable must undergo
modification $\out{\delta_{v}}$,

\item \textbf{Bottom right}: and the variable's type has undergone modification
$\out{\delta_{\tau}}$.

\end{itemize}

The rules in Figure~\ref{lookup-diff-dispatch} summarize the global strategy of
first looking up in the local context, and then falling back to the global
environment, analogously to the process in Figure~\ref{lookup-type-dispatch}.

\begin{Rules}
  {lookup-diff-dispatch}
  { Diff lookup rules (local context and global environment) }

\begin{mathpar}
  {
    \inferrule*
    [lab=Diff-\LookupCtxt]
    {\turnstile%
      { \diff{\Gamma}{\delta_\Gamma} }
      { \declDiff{ v }{ \delta_v }{ \delta_{\tau} } }
    }
    {\turnstile%
      {\dcontext{E}{\delta_E}{\Gamma}{\delta_\Gamma} }
      { \declDiff{ v }{ \delta_v }{ \delta_{\tau} } }
    }
  }

  {
    \inferrule*
    [lab=Diff-\LookupEnv]
    {
      {\turnstile%
        { \diff{E} {\delta_E} }
        { \declDiff{ v } { \delta_v } { \delta_{\tau} } }
      }
    }
    {\turnstile%
      { \dcontext{E} {\delta_E} {\Gamma} {\delta_\Gamma} }
      { \declDiff{ v } { \delta_v } { \delta_{\tau} } }
    }
  }

\end{mathpar}

\end{Rules}

\subsubsection{Lookup rules for diffs in the local context}

Figure~\ref{lookup-diff-context} shows how we proceed to lookup for the diff of
a variable from the original program in the local context.  Note that, since we
assume the original program type-checked, when we look up variables, we can
assume that they are properly bound: were it not the case, we would need
additional rules to account for the absence of a variable, and to change the
output types to allow for failure.

When the context hasn't changed, we can obtain its type from the old environment
(Rule~\rulename{Diff-\LCSame}).  When some variable has been inserted, we can
ignore it with no effort, since we only care about shadowing in the original
program (Rule~\rulename{Diff-\LCIns}).

When we find the target variable in the local context, there are only two
remaining possibilities for the diff: either it is a $\MathModPiOp$, or a
$\MathDropPiOp$.  When it has been modified, we have found exactly the
information we were looking for, so we can simply output it
(Rule~\rulename{Diff-\LCMod-Here}).  If it has been dropped, then there is no
satisfactory result, the variable is no longer in the program.  In order to make
sure the user notices, we rename the variable using the helper
$\op{\text{Deprecate}}$, and make its diff as uninformative as possible using a
hole (Rule~\rulename{Diff-\LCMod-There}).

The last cases to handle are that of a variable other than our target being
modified (Rule~\rulename{Diff-\LCDrop}), in which case we simply keep looking
recursively into the context, and that of a permutation
(Rule~\rulename{Diff-\LCPermute}), in which case we can perform the permutation
and keep looking recursively.

\begin{Rules}{lookup-diff-context}{ Diff lookup rules (local context) }

  \begin{mathpar}
    {
      \inferrule*
      [lab=Diff-\LCSame]
      {
        { \turnstile {\Gamma} {\hasType { v } { \out{\tau_v} } } }
      }
      {\turnstile
        { \diff {\Gamma} {\MathSame} }
        { \declDiff { v } { \MathSame } { \MathSame } }
      }
    }

    {
      \inferrule*
      [lab=Diff-\LCIns]
      {\turnstile
        { \diff {\Gamma} {\delta_{\Gamma}} }
        { \declDiff { v } { \delta_v } { \delta_{\tau_v} } }
      }
      {\turnstile
        { \diff {\Gamma} {\MathIns{ (\delta_w : \delta_{\tau_{w}}) }{\delta_{\Gamma}}} }
        { \declDiff { v } { \delta_v } { \delta_{\tau_v} } }
      }
    }

    {
      \inferrule*
      [lab=Diff-\LCMod-Here]
      {
      }
      {\turnstile
        { \diff {\MathCons{\MathLocalAssum{v}{\tau}}{\Gamma}} {\MathMod{(\delta_v : \delta_{\tau})}{\delta_{\Gamma}}} }
        { \declDiff { v } { \delta_v } { \delta_{\tau} } }
      }
    }

    {
      \inferrule*
      [lab=Diff-\LCMod-There]
      {
        w \neq v
        \and
        {\turnstile
          { \diff {\Gamma} {\delta_{\Gamma}} }
          { \declDiff { v } { \delta_v } { \delta_{\tau} } }
        }
      }
      {\turnstile
        { \diff {\MathCons{\MathLocalAssum{w}{\tau}}{\Gamma}} {\MathMod{(\delta_v : \delta_{\tau})}{\delta_{\Gamma}}} }
        { \declDiff { v } { \delta_v } { \delta_{\tau} } }
      }
    }

    {
      \inferrule*
      [lab=Diff-\LCDrop]
      { \op{\text{Deprecate}(} v \op{)} = \out{v'} }
      {\turnstile%
        { \diff%
          { \MathCons{\MathLocalAssum{v}{\tau}}{\Gamma} }
          { \MathDrop{\delta_{\Gamma}} }
        }
        { \declDiff { v } { \MathReplace{v'} } { \MathReplace{\MathHole{}} } }
      }
    }

    {
      \inferrule*
      [lab=Diff-\LCPermute]
      {\turnstile%
        { \diff%
          { \permuteOp{p}{\Gamma} }
          { \delta_{\Gamma} }
        }
        { \declDiff { v } { \delta_{v} } { \delta_{\tau} } }
      }
      {\turnstile%
        { \diff%
          { \Gamma }
          { \MathPermute{p}{\delta_{\Gamma}} }
        }
        { \declDiff { v } { \delta_{v} } { \delta_{\tau} } }
      }
    }

\end{mathpar}

\end{Rules}

\subsubsection{Lookup rules for diffs in the global environment}

\begin{Rules}
{rules-lookup-environment}
{Lookup in the global environment}

  \begin{mathpar}

    {
      \inferrule*
      [lab=Diff-\GESame]
      {
        { \turnstile {E} {\hasType { v } { \out{\tau_v} } } }
      }
      {\turnstile
        { \diff {E} {\MathSame} }
        { \declDiff { v } { \MathSame } { \MathSame } }
      }
    }

    {
      \inferrule*
      [lab=Diff-\GEDrop-Definition-Here]
      {
        {
          \op{\text{Deprecate}(} v \op{)} = \out{v'}
        }
      }
      {
        \turnstile%
        { \diff%
          {
            \MathCons%
            {\Definition{k}{v}{\tau}{t}}
            {E}
          }
          {\MathDrop{\delta_{E}}} }
        { \declDiff{v}{\MathReplace{v'}}{\MathReplace{\MathHole{}}} }
      }
    }

    {
      \inferrule*
      [lab=Diff-\GEDrop-Definition-There]
      {
        { w \neq v }
        \and
        {
          \turnstile%
          { \diff%
            {E}
            {\delta_{E}}
          }
          { \declDiff{v}{\delta_v}{\delta_\tau} }
        }
      }
      {
        \turnstile%
        { \diff%
          {
            \MathCons%
            {\Definition{k}{w}{\tau}{t}}
            {E}
          }
          {\MathDrop{\delta_{E}}} }
        { \declDiff{v}{\delta_v}{\delta_\tau} }
      }
    }

    {
      \inferrule*
      [lab=Diff-\GEIns]
      {\turnstile
        { \diff {E} {\delta_{E}} }
        { \declDiff { v } { \delta_v } { \delta_{\tau_v} } }
      }
      {\turnstile
        { \diff%
          {E}
          {
            \MathIns{
              \Inductive%
              {n}
              {\overline{p}}
              {\overline{i}}
              {u}
              {\overline{c}}
            }
            {\delta_{E}}
          }
        }
        { \declDiff { v } { \delta_v } { \delta_{\tau_v} } }
      }
    }

    {
      \inferrule*
      [lab=Diff-\GEMod{Defn}]
      {
      }
      {\turnstile
        { \diff
          {\MathCons{\Definition{k}{v}{\tau}{t}}{E}}
          {\MathMod{\ModifyDefinition{\delta_k}{\delta_v}{\delta_\tau}{\delta_t}}{\delta_{E}}}
        }
        { \declDiff { v } { \delta_v } { \delta_{\tau} } }
      }
    }

    {
      \inferrule*
      [lab=Diff-\GEMod{Inductive}-Type]
      {
        {
          I = \Inductive%
          {n}
          {\overline{p}}
          {\overline{i}}
          {u}
          {\overline{c}}
        }
        \\
        {
          \delta_{I} = \ModifyInductive%
          {\delta_n}
          {\delta_{\overline{p}}}
          {\delta_{\overline{i}}}
          {\delta_u}
          {\delta_{\overline{c}}}
        }
        \\\\
        {
          \deltaIndType{I}{\delta_{I}}{\delta_\tau}
        }
      }
      {
        \turnstile%
        {
          \diff%
          {\MathCons{I}{E}}
          {
            \MathMod%
            {\delta_{I}}
            {\delta_{E}}
          }
        }
        { \declDiff{n}{\delta_n}{\delta_{\tau}} }
      }
    }

    % {
    %   \inferrule*
    %   [lab=Diff-\GEMod{Inductive}]
    %   { T
    %   }
    %   { B
    %   }
    % }

    {
      \inferrule*
      [lab=Diff-\GEMod{Eliminator}]
      {
        {
          I = \Inductive%
          {n}
          {\overline{p}}
          {\overline{i}}
          {u}
          {\overline{c}}
        }
        \\
        {
          \delta_{I} = \ModifyInductive%
          {\delta_n}
          {\delta_{\overline{p}}}
          {\delta_{\overline{i}}}
          {\delta_u}
          {\delta_{\overline{c}}}
        }
        \\\\
        {
          v = \op{\text{EliminatorName}(} I \op{,\ } u_{\text{elim}} \op{)}
        }
        \and
        {
          \deltaEliminatorName{I}{\delta_{I}}{u_{\text{elim}}}
          {\delta_v}
        }
        \\
        {
          \deltaEliminator{I}{\delta_{I}}{u_{\text{elim}}}
          {\delta_\tau}
        }
      }
      {\turnstile%
        { \diff%
          {\MathCons{I}{E}}
          {\MathMod{\delta_{I}}{\delta_{E}}}
        }
        { \declDiff{v}{\delta_v}{\delta_{\tau}} }
      }
    }

  \end{mathpar}

\end{Rules}

Unfortunately, the presentation for $\ConstructorText{}$ is quite heavy, as we
need to account for all possible list changes to the list of constructors.  In
the interest of thoroughness, we list all the rules in
Figure~\ref{rules-lookup-environment-constructor}, but do warn the reader that
reading these rules is not particularly interesting or enlightening.

\begin{Rules}
{rules-lookup-environment-constructor}
{Lookup in the global environment (constructor rules)}

  \begin{mathpar}

    % Possibilities:
    % (v + cs), (δv +mod δcs)
    % (v + cs), (-drop δcs)
    % (v + cs), (+keep δcs)
    % (v + cs), (permute p δcs)
    % (v + cs), (v' +ins δcs)
    % (v' + cs), (δv +mod δcs)
    % (v' + cs), (-drop δcs)
    % (v' + cs), (+keep δcs)
    % (v' + cs), (permute p δcs)
    % (v' + cs), (v'' +ins δcs)

    {
      \inferrule*
      [lab=Diff-\GEMod{Constructor}]
      {
        {
          I = \Inductive%
          {n}
          {\overline{p}}
          {\overline{i}}
          {u}
          {\MathCons{C}{\overline{c}}}
        }
        \\
        {
          \delta_{I} = \ModifyInductive%
          {\delta_n}
          {\delta_{\overline{p}}}
          {\delta_{\overline{i}}}
          {\delta_u}
          {\delta_{\overline{c}}}
        }
        \\
        {
          C = \Constructor%
          {n_{\text{ctor}}}
          {\overline{p_{\text{ctor}}}}
          {\overline{i_{\text{ctor}}}}
        }
        \\
        {
          \deltaConstructor
          {\overline{p}}
          {\delta_{\overline{p}}}
          {\overline{i}}
          {\delta_{\overline{i}}}
          {\overline{c}}
          {\delta_{\overline{c}}}
          {v}
          {(\delta_v, \delta_\tau)}
        }
      }
      {\turnstile
        { \diff
          {\MathCons
            {\Inductive
              {n}
              {\overline{p}}
              {\overline{i}}
              {u}
              {\overline{c}}
            }
            {E}
          }
          {\MathMod
            {\ModifyInductive
              {\delta_n}
              {\delta_{\overline{p}}}
              {\delta_{\overline{i}}}
              {\delta_u}
              {\delta_{\overline{c}}}
            }
            {\delta_{E}}
          }
        }
        { \declDiff{v}{\delta_v}{\delta_{\tau}} }
      }
    }

  \end{mathpar}

\end{Rules}

\documentclass[preview]{standalone}

\usepackage{fontspec}
\usepackage{gfsdidot}          % makes math stuff look neat, but also overrides text stuff
\setmainfont{TeX Gyre Pagella} % forces back text stuff to look neat too
\setmonofont{DejaVuSansMono}   % monospace stuff should look neat too

\usepackage[dvipsnames]{xcolor} % must be loaded prior to tikz, pgfplots, etc.
\let\iint\relax
\let\iiint\relax
\let\iiiint\relax
\let\idotsint\relax
\usepackage{amsmath} % loaded by mathtools?
\usepackage{amssymb}
\usepackage[english]{babel}
\usepackage{booktabs}
\usepackage{calc}
\usepackage{caption}
\usepackage{color}
\usepackage{DejaVuSansMono}
\usepackage{dsfont}
\let\textlozenge\relax
\usepackage[inline]{enumitem}
\usepackage{environ}
\usepackage{etoolbox}
\usepackage{fontawesome}
\usepackage{forest}
\usepackage{graphicx}
\usepackage{longtable}
\usepackage{ltablex} \keepXColumns % Made me sad for appendix tables, where was it used?
\usepackage{makecell}
\usepackage{mathpartir}
%\usepackage{mathtools} % \coloneqq
\usepackage{mdframed}
\usepackage{minted}
\usepackage{multicol}
\usepackage[parfill]{parskip}
\usepackage{pgfplots}
\usepackage{pifont} % \ding
\usepackage{soul}
\usepackage{syntax}
\usepackage{tabularx}
\usepackage[skins,theorems]{tcolorbox}
\usepackage{tikz}
\usepackage{tikzpeople}
\usepackage{tikz-qtree}
\usepackage{xpatch}
\usepackage[breaklinks=true,pdfborder={0 0 0}]{hyperref}

% gfsdidot makes ∀ and ∃ super ugly, this reverts them to something fine
\DeclareSymbolFont{CMsymbols}{OMS}{cmsy}{b}{n}
\SetSymbolFont{CMsymbols}{bold}{OMS}{cmsy}{ub}{n}
\DeclareMathSymbol{\forall}{\mathord}{CMsymbols}{"38}
\DeclareMathSymbol{\exists}{\mathord}{CMsymbols}{"39}

\definecolor{monokaibg}{HTML}{272822}
\definecolor{color01}{HTML}{E6194B}
\definecolor{color02}{HTML}{3CB44B}
\definecolor{color03}{HTML}{FFE119}
\definecolor{color04}{HTML}{4363D8}
\definecolor{color05}{HTML}{F58231}
\definecolor{color06}{HTML}{911EB4}
\definecolor{color07}{HTML}{46F0F0}
\definecolor{color08}{HTML}{F032E6}
\definecolor{color09}{HTML}{BCF60C}
\definecolor{color10}{HTML}{FABEBE}
\definecolor{color11}{HTML}{008080}
\definecolor{color12}{HTML}{E6BEFF}
\definecolor{color13}{HTML}{9A6324}
\definecolor{color14}{HTML}{FFFAC8}
\definecolor{color15}{HTML}{800000}
\definecolor{color16}{HTML}{AAFFC3}
\definecolor{color17}{HTML}{808000}
\definecolor{color18}{HTML}{FFD8B1}
\definecolor{color19}{HTML}{000075}
\definecolor{color20}{HTML}{808080}
\definecolor{color21}{HTML}{FFFFFF}
\definecolor{color22}{HTML}{000000}

\def\mathunderline#1#2{\color{#1}\underline{{\color{black}#2}}\color{black}}

\forestset{
  default preamble={
    for tree={
      align=center,
      draw=black,
      edge={
        ultra thick,
      },
      edge path={
        \noexpand\path[\forestoption{edge}]
        (!u.parent anchor)
        -- +(0,-10pt)
        -| (.child anchor)
        \forestoption{edge label};
      },
      font=\bfseries,
      % inner sep=10pt,
      l sep=30pt,
      % line width=2pt,
      minimum height=20pt,
      minimum width=30pt,
      parent anchor=south,
      rectangle,
      ultra thick,
    }
  }
}

\newcolumntype{Y}{>{\centering\arraybackslash}X}

% defines safecoqinline, a command like coqinline but works in tabular
% environments
\makeatletter
\def\safecoqinline#1{%
\ifx\@footnotetext\TX@trial@ftn
\detokenize{#1}%
\else
\coqinline{#1}%
\fi}
\makeatother

\captionsetup{
  justification=centering,
}

\usetikzlibrary{arrows, calc, fit, positioning, shapes, tikzmark}

\tikzset{
bicolor/.style 2 args={
  dashed,dash pattern=on 4pt off 4pt,#1,
  postaction={draw,dashed,dash pattern=on 4pt off 4pt,#2,dash phase=4pt}
  },
}

\tikzstyle{Matching}=[
black,
dashed,
line width=3pt,
]

\tikzstyle{RoundedDottedPath}=[
densely dotted,
color05,
line width=3pt,
rounded corners=10pt,
]

\tikzstyle{RoundedRectangle}=[
dashed,
draw,
line width=3pt,
rounded corners=10pt,
]

\tikzstyle{NodeLabel}=[
black,
circle,
draw,
fill=white,
font=\bfseries,
inner sep=1pt,
ultra thick,
]

\graphicspath{ {./images/} }

% Fix a mdframed bug where skipbelow is ignored
% \makeatletter
% \xpatchcmd{\endmdframed}
%   {\aftergroup\endmdf@trivlist\color@endgroup}
%   {\endmdf@trivlist\color@endgroup\@doendpe}
%   {}{}
% \makeatother

% \surroundwithmdframed[
% backgroundcolor=monokaibg,
% linecolor=white,
% skipabove=1em,
% skipbelow=1em,
% leftmargin=10pt,
% innertopmargin=1pt,
% innerbottommargin=0pt
% ]{minted}

\setminted{
  % DO NOT use bgcolor, use mdframed instead
  % bgcolor=monokaibg,
  linenos=true,
  style=lovelace,
  breaklines=true,
  encoding=utf8,
  fontsize=\large,
  baselinestretch=1,
}

% need bgcolor for the display to be nice in rules
\newmintinline{coq}{bgcolor=white}

\makeatletter
\pgfdeclareshape{document}{
\inheritsavedanchors[from=rectangle] % this is nearly a rectangle
\inheritanchorborder[from=rectangle]
\inheritanchor[from=rectangle]{center}
\inheritanchor[from=rectangle]{north}
\inheritanchor[from=rectangle]{south}
\inheritanchor[from=rectangle]{west}
\inheritanchor[from=rectangle]{east}
% ... and possibly more
\backgroundpath{% this is new
% store lower right in xa/ya and upper right in xb/yb
\southwest \pgf@xa=\pgf@x \pgf@ya=\pgf@y
\northeast \pgf@xb=\pgf@x \pgf@yb=\pgf@y
% compute corner of ‘‘flipped page’’
\pgf@xc=\pgf@xb \advance\pgf@xc by-10pt % this should be a parameter
\pgf@yc=\pgf@yb \advance\pgf@yc by-10pt
% construct main path
\pgfpathmoveto{\pgfpoint{\pgf@xa}{\pgf@ya}}
\pgfpathlineto{\pgfpoint{\pgf@xa}{\pgf@yb}}
\pgfpathlineto{\pgfpoint{\pgf@xc}{\pgf@yb}}
\pgfpathlineto{\pgfpoint{\pgf@xb}{\pgf@yc}}
\pgfpathlineto{\pgfpoint{\pgf@xb}{\pgf@ya}}
\pgfpathclose
% add little corner
\pgfpathmoveto{\pgfpoint{\pgf@xc}{\pgf@yb}}
\pgfpathlineto{\pgfpoint{\pgf@xc}{\pgf@yc}}
\pgfpathlineto{\pgfpoint{\pgf@xb}{\pgf@yc}}
\pgfpathlineto{\pgfpoint{\pgf@xc}{\pgf@yc}}
}
}
\makeatother

%%%%% General-purpose text macros %%%%%
\newcommand{\define}[1]{\emph{#1}}
\newcommand{\mycite}[1]{\citeauthor{#1}~\cite{#1}}

\newcommand{\Language}[1]{\emph{#1}}
\newcommand{\Tool}[1]{\emph{#1}}
\newcommand{\Gallina}{\Language{Gallina}}
\newcommand{\Ltac}{\Language{Ltac}}
\newcommand{\Vernacular}{\Language{Vernacular}}

\newcommand{\Chick}{\Language{Chick}}
\newcommand{\Coq}{\Language{Coq}}
\newcommand{\CoqIDE}{\Tool{CoqIDE}}
\newcommand{\Haskell}{\Language{Haskell}}
\newcommand{\JavaScript}{\Language{JavaScript}}
\newcommand{\OCaml}{\Language{OCaml}}
\newcommand{\PeaCoq}{\Language{PeaCoq}}
\newcommand{\RxJS}{\Language{RxJS}}
\newcommand{\SerAPI}{\Tool{SerAPI}}
\newcommand{\Snap}{\Language{Snap}}
\newcommand{\TypeScript}{\Language{TypeScript}}

\newcommand{\OperatorColor}{purple}
\newcommand{\Operator}[1]{\textcolor{\OperatorColor}{\ #1\ }}
\newcommand{\Entails}{\Operator{\vdash}}
\newcommand{\HasType}{\Operator{:}}

%\newmintinline{coq}{fontsize=\small}
% \newcommand{\coqinline}[1]{%
%   %\colorbox{monokaibg}{%
%   \parbox[c][0.9em]{\widthof{\mycoq{#1}}}{\mycoq{#1}}%
%   %}%
% }

\newcommand{\modified}[1]{\color{Orange}{#1}}
\newcommand{\repaired}[1]{\color{PineGreen}{#1}}

\newcommand{\rulename}[1]{$\LeftTirNameStyle{#1}$}

\newcommand{\RmApp}{Rm-App}
\newcommand{\RmPi}{Rm-Pi}
\newcommand{\InsApp}{Ins-App}

%%%%% Math-mode macros %%%%%

\newcommand{\opcolor}{purple}
\newcommand{\out}[1]{ \boxed{ \textcolor{teal}{#1} } }
\newcommand{\MathPatches}[3]{%
#1 \overset{#2}{\mathbin{\textcolor{\opcolor}{\rightsquigarrow}}} \out{#3}%
}

\newcommand{\App}{\$}
\newcommand{\Mod}{Mod}
\newcommand{\Drop}{Drop}
\newcommand{\Ins}{Ins}
\newcommand{\Keep}{Keep}
\newcommand{\Lam}{\lambda}

\newcommand{\oMod}{\overset{\mathtt{\Mod}}}
\newcommand{\oDrop}{\overset{\mathtt{\Drop}}}
\newcommand{\oIns}{\overset{\mathtt{\Ins}}}
\newcommand{\oKeep}{\overset{\mathtt{\Keep}}}

\newcommand{\Cons}{::}
\newcommand{\permute}[2]{%
\overline{#2}^{\overset{#1}{\rightleftarrows}}%
}
\newcommand{\permuteOp}[2]{%
\overset{\overset{#1}{\rightleftarrows}}{#2}%
}

\newcommand{\mkMathPiRaw}[4]{#3{\Pi} #4{#1} \rightarrow #2}
\newcommand{\mkMathPi}[5]{\mkMathPiRaw{(#2 : #1)}{#3}{#4}{#5}}


\newcommand{\MathCons}[2]{#1 \Cons #2}
\newcommand{\MathDrop}[1]{\oDrop{\Cons} #1}
\newcommand{\MathDropApp}[1]{\oDrop{\App} #1}
\newcommand{\MathDropLam}[1]{\oDrop{\Lam} #1}
\newcommand{\MathDropPi}[1]{\oDrop{\Pi} #1}
\newcommand{\MathHole}{\texttt{\_}}
\newcommand{\MathIns}[2]{#1 \oIns{\Cons} #2}
\newcommand{\MathInsApp}[2]{\mkMathApp{#1}{#2}{\oIns}}
\newcommand{\MathInsAppOp}{\oIns{\App}}
\newcommand{\MathInsLam}[2]{\mkMathLam{#1}{#2}{\oIns}{}}
\newcommand{\MathInsLamOp}{\oIns{\Lam}}
\newcommand{\MathInsPi}[3]{\mkMathPi {#1}{#2}{#3}{\oIns}{}}
\newcommand{\MathInsPiOp}{\oIns{\Pi}}
\newcommand{\MathLam}[2]{\mkMathLam{#1}{#2}{}{}}
\newcommand{\MathLams}[3]{\mkMathLam{#1}{#2}{}{\BarCount{#3}}}
\newcommand{\MathKeepPiOp}{\oKeep{\Pi}}
\newcommand{\MathMod}[2]{#1 \oMod{\Cons} #2}
\newcommand{\MathModAppOp}{\oMod{\App}}
\newcommand{\MathModLamOp}{\oMod{\Lam}}
\newcommand{\MathModPiOp}{\oMod{\Pi}}
\newcommand{\MathPermute}[2]{\permuteOp{#1}{\Cons} #2}
\newcommand{\MathPermuteLams}[2]{\permuteOp{#1}{\lambda} #2}
\newcommand{\MathPermutePis}[2]{\permuteOp{#1}{\Pi} #2}
\newcommand{\MathPi}[3]{\mkMathPi{#1}{#2}{#3}{}{}}
\newcommand{\MathPis}[4]{\mkMathPi{#1}{#2}{#3}{}{\BarCount{#4}}}
\newcommand{\MathReplace}[1]{\mathds{K}({#1})}
\newcommand{\MathSame}{\mathds{1}}

\newcommand{\mkMathApp}[3]{#1 #3{ \$ } #2}

%%% REPAIR %%%

\newcommand{\op}[1]{\textcolor{\opcolor}{#1}}

\newcommand{\boundOp}{\op{\text{Bound}}}
\newcommand{\freshOneOp}{\op{\text{Fresh}_1}}
\newcommand{\freshTwoOp}{\op{\text{Fresh}_2}}
\newcommand{\repairIndOp}{\op{R_{I}}}
\newcommand{\repairProgOp}{\op{R_{P}}}
\newcommand{\repairTermOneOp}{\op{R_{T_1}}}
\newcommand{\repairTermTwoOp}{\op{R_{T_2}}}
\newcommand{\repairTermThreeOp}{\op{R_{T_2}}}
\newcommand{\repairVernacOp}{\op{R_{V}}}
\newcommand{\vernacEnvOp}{\op{E_{V}}}
\newcommand{\repairBranchesOp}{\op{R_{B}}}

\newcommand{\blackbrackets}[1]{\left[ #1 \right]}
\newcommand{\blackdiff}[2]{\blackbrackets{\genfrac{}{}{0pt}{}{#1}{#2}}}
\newcommand{\bound}[1]{\boundOp\op{(}#1\op{)}}
\newcommand{\brackets}[1]{%
  \color{\opcolor} \left[ \normalcolor #1 \color{\opcolor} \right] \normalcolor%
}
\newcommand{\context}[2]{#1 \op{,} #2}
\newcommand{\dcontext}[4]{\context{\diff{#1}{#2}}{\diff{#3}{#4}}}
\newcommand{\denv}[2]{\diff{#1}{#2}}
\newcommand{\diff}[2]{\brackets{\genfrac{}{}{0pt}{}{#1}{#2}}}
\newcommand{\dtau}[1]{\delta_{\tau_{#1}}}
\newcommand{\equals}[2]{#1 \mathbin{\textcolor{\opcolor}{=}} #2}
\newcommand{\freshOne}[3]{\freshOneOp\op{(}\diff{#1}{#2}\op{) =}\ \out{#3}}
\newcommand{\freshTwo}[3]{\freshTwoOp\op{(}#1\op{,}#2\op{) =}\ \out{#3}}
\newcommand{\genericrepair}[3]{\mkrepair{\repairTermTwoOp}{#1}{#2}{#3}{\op{?}}}
\newcommand{\nturnstile}[2]{#1 \mathbin{\textcolor{\opcolor}{\nvdash}} #2}
\newcommand{\qmark}{\textcolor{\opcolor}{?}}
\newcommand{\squiggly}[2]{#1 \mathbin{\textcolor{\opcolor}{\rightsquigarrow}} #2}
\newcommand{\subst}[2]{#1 \leftarrow #2}
\newcommand{\turnstile}[2]{#1\ \mathbin{\textcolor{\opcolor}{\vdash}}\ #2}
\newcommand{\repair}[4]{ \mkrepair{\repairTermOneOp}{#1}{#2}{#3}{#4} }
\newcommand{\repairBranches}[3]{\mksimplerepair{\repairBranchesOp}{#1\op{,}\ #2}{#3}}
\newcommand{\repairInd}[7]{
  \mksimplerepair%
  {\repairIndOp}
  {\diff{#1}{#2}\op{,}\ \diff{#3}{#4}\op{,}\ \diff{#5}{#6}}
  {#7}
}
\newcommand{\repairProg}[3]{
  \repairProgOp\op{(}\blackdiff{#1}{#2}\op{)}\ \op{=}\ \out{#3}
}
\newcommand{\repairTerm}[4]{\mkrepair{\repairTermOneOp}{#1}{#2}{#3}{#4}}
\newcommand{\repairTermWithoutType}[2]{\mksimplerepair{\repairTermThreeOp}{#1}{#2}}
\newcommand{\repairVernac}[3]{
  \repairVernacOp\op{(}\blackdiff{#1}{#2}\op{)}\ \op{=}\ \out{#3}
}

\newcommand{\DefinitionText}{\text{Definition}}
\newcommand{\InductiveText}{\text{Inductive}}
\newcommand{\ConstructorText}{\text{Constructor}}
\newcommand{\GlobalDefinitionText}{\text{GlobalDefinition}}
\newcommand{\GlobalInductiveText}{\text{GlobalInductive}}

\newcommand{\Definition}[4]{ \DefinitionText(\{ #1, #2, #3, #4 \}) }
\newcommand{\Inductive}[5]{ \InductiveText(\{ #1, #2, #3, #4, #5 \}) }
\newcommand{\Constructor}[3]{ \ConstructorText(\{ #1, #2, #3 \}) }
\newcommand{\GlobalDefinition}[3]{ \GlobalDefinitionText(\{ #1, #2, #3 \}) }
\newcommand{\GlobalInductive}[5]{ \GlobalInductiveText(\{ #1, #2, #3, #4, #5 \}) }

\newcommand{\ModifyDefinition}[4]{ \delta_{\DefinitionText}(\{ #1, #2, #3, #4 \}) }
\newcommand{\ModifyInductive}[5]{ \delta_{\InductiveText}(\{ #1, #2, #3, #4, #5 \}) }
\newcommand{\ModifyGlobalDefinition}[3]{ \delta_{\GlobalDefinitionText}(\{ #1, #2, #3 \}) }
\newcommand{\ModifyGlobalInductive}[5]{ \delta_{\GlobalInductiveText}(\{ #1, #2, #3, #4, #5 \}) }

\newcommand{\GlobalDefinitionAnon}{\GlobalDefinitionText(\{ \ldots \})}
\newcommand{\GlobalInductiveAnon}{\GlobalInductiveText(\{ \ldots \})}

\newcommand{\DefinitionAnon}{\DefinitionText(\{ \ldots \})}
\newcommand{\InductiveAnon}{\InductiveText(\{ \ldots \})}

\newcommand{\RepairProg}[1]{R-Prog-#1}

\newcommand{\RProgMod}{\RepairProg{Modify}}
\newcommand{\RProgKeep}{\RepairProg{Keep}}
\newcommand{\RProgSameCons}{\RepairProg{Same-Cons}}

\newcommand{\scalefactor}{1}
\newcommand{\invscalefactor}{1}

\NewEnviron
    {Rules}[2]
    {
      \begin{figure*}
        %\scalebox{\invscalefactor}{
          %\begin{minipage}{\scalefactor\textwidth}
            %\begin{mathpar}
              \BODY
            %\end{mathpar}
          %\end{minipage}
        %}
        \caption{#2}
        \label{#1}
      \end{figure*}
    }

\newcommand{\MathSameProg}{\MathSame_P}
\newcommand{\MathSameVernac}{\MathSame_V}
\newcommand{\MathSameGlobal}{\MathSame_G}

\newcommand{\MathProp}{\mathtt{Prop}}
\newcommand{\MathSet} {\mathtt{Set}}
\newcommand{\MathType}{\mathtt{Type}}

\newcommand{\hasType}[2]{ #1\ \op{:}\ #2 }

\newcommand{\mkrepair}[5]{
  #1
  \op{(}
  \hasType
  { #2 }
  { \diff{ #4 }{ #5 } }
  \op{) =\ }
  \out{ #3 }
}

\newcommand{\MathLocalAssum}[2]{ ( #1 : #2 ) }
\newcommand{\MathLocalDef}[3]{ ( #1 : #2 \coloneqq #3 ) }

\newcommand{\mksimplerepair}[3]{
  #1
  \op{(}
  #2
  \op{) =\ }
  \out{ #3 }
}

\newcommand{\RModPi}{\RepairTerm{Mod-$\Pi$}}
\newcommand{\RInsPi}{\RepairTerm{Ins-$\Pi$}}
\newcommand{\RDropPi}{\RepairTerm{Drop-$\Pi$}}
\newcommand{\RPermutePis}{\RepairTerm{Permute-$\Pi$s}}
\newcommand{\RReplace}{\RepairTerm{Replace}}
\newcommand{\RSame}{\RepairTerm{$\MathSame{}$-Other}}
\newcommand{\RSamePi}{\RepairTerm{$\MathSame{}$-$\Pi$}}

\newcommand{\RepairTermPrefix}{RT1}
\newcommand{\GenericRepairPrefix}{R-Term-2}
\newcommand{\UnknownTypeRepairPrefix}{R-Term-2}

\newcommand{\RepairTerm}[1]{\RepairTermPrefix-#1}
\newcommand{\GenericRepair}[1]{\GenericRepairPrefix-#1}
\newcommand{\UnknownTypeRepair}[1]{\UnknownTypeRepairPrefix-#1}

\newcommand{\mkMathLam}[4]{#3{\lambda} #4{#1} \rightarrow #2}

\newcommand{\MathModApp}[2]{ \mkMathApp{#1}{#2}{\oMod}{} }
\newcommand{\MathModLam}[2]{ \mkMathLam{#1}{#2}{\oMod}{} }
\newcommand{\MathModPi} [3]{ \mkMathPi{#1}{#2}{#3}{\oMod}{} }
\newcommand{\MathKeepPi}[1]{ \MathKeepPiOp \rightarrow #1 }

\newcommand{\MathAnnot}[2]{#1 : #2}

\newcommand{\UTRApp}{\UnknownTypeRepair{App}}
\newcommand{\UTRMatch}{\UnknownTypeRepair{Match}}
\newcommand{\UTRPi}{\UnknownTypeRepair{Pi}}
\newcommand{\UTRType}{\UnknownTypeRepair{Universe}}
\newcommand{\UTRVar}{\UnknownTypeRepair{Var}}
\newcommand{\UTRAnnot}{\UnknownTypeRepair{Annot}}
\newcommand{\UTRHole}{\UnknownTypeRepair{Hole}}
\newcommand{\UTROtherwise}{\UnknownTypeRepair{Otherwise}}

\newcommand{\repairArgsOp}{\op{R_{Fn}}}

\newcommand{\repairArgs}[5]{
  \op{\repairArgsOp(}#1\op{,} #2\op{,} #3\op{,} #4\op{) =}\ \out{#5}
}

\newcommand{\declDiff}[3]{
  \hasType
      { \diff{ #1 }{ \out{#2} } }
      { \diff{ \op{?} }{ \out{#3} } }
}

\newcommand{\MathMatch}[2]{ \text{match}\ #1\ \text{with}\ #2 }


\begin{document}


\newcommand{\RepairTermOneSamePi}{
  \inferrule*
  [lab=\RSamePi]
  {
    {\turnstile
      { \dcontext{E}{\delta_E}{\Gamma}{\delta_\Gamma} }
      { \repair
        {t}
        {\delta}
        {\MathPi{\tau_1}{\chi}{\tau_2}}
        {\MathModPi{\MathSameTerm{}}{\MathSameBinder{}}{\MathSameTerm{}}}
      }
    }
  }
  {\turnstile
    { \dcontext{E}{\delta_E}{\Gamma}{\delta_\Gamma} }
    { \repair
      {t}
      {\delta}
      {\MathPi{\tau_1}{\chi}{\tau_2}}
      {\MathSameTerm{}}
    }
  }
}

\newcommand{\RepairTermOneSameOther}{
  \inferrule*
  [lab=\RSame]
  {
    \text{\rulename{\RSamePi} does not apply}
    \and
    {\turnstile
      { \dcontext{E}{\delta_E}{\Gamma}{\delta_\Gamma} }
      { \repairTermWithoutType{t}{\delta_{t}}}
    }
  }
  {\turnstile
    { \dcontext{E}{\delta_E}{\Gamma}{\delta_\Gamma} }
    { \repair{t}{\delta_{t}}{\tau}{\MathSame} }
  }
}

\newcommand{\RepairTermOneModPi}{
  \inferrule*
  [lab=\RModPi]
  {
    {\turnstile%
      {\dcontext{E}{\delta_E}{\Gamma}{\delta_\Gamma}}
      {\freshOne{x}{\MathSame}{\delta_x}}
    }
    \and
    {\MathPatches{x}{\delta_x}{x'}}
    % \and
    % {\diffExpand{\tau_{2}}{\delta_{\tau_{2}}}{\delta_{\tau_{2}}'}
    % }
    \\\\\\
    {\turnstile%
      {\dcontext{E}{\delta_E}{\Gamma}{\delta_\Gamma}}
      {\freshTwo{\blackdiff{x}{\delta_x}}{\blackdiff{\chi}{\delta_\chi}}{z}}
    }
    \and
    {\MathPatches{\chi}{\delta_\chi}{\chi'}}
    \\\\\\
    {\turnstile%
      { \dcontext{E}{\delta_E}
        { \MathCons{\MathLocalAssum{z}{\tau_1}}{\Gamma} }
        { \MathMod%
          {\MathLocalAssum{\MathReplace{z}}{\delta_{\tau_1}}}
          {\delta_\Gamma}
        }
      }
      { \repair%
        {t[\subst{x}{z}]}
        {\delta_{t}}
        {\tau_2[\subst{\chi}{z}]}
        {\delta_{\tau_2}[\subst{\MathReplace{\chi'}}{\MathReplace{z}}]}
      }
    }
    % \\\\\\
    % {\diffExpand{t}{\delta_{t}}{\delta_{t}'}
    % }
  }
  {\turnstile%
    { \dcontext{E}{\delta_E}{\Gamma}{\delta_\Gamma} }
    {\repair%
      { \MathLam{x}{t} }
      { \MathModLam{\delta_x}{\delta_{t}[\subst{\MathReplace{z}}{\MathReplace{x'}}]} }
      { \MathPi{\tau_1}{\chi}{\tau_2} }
      { \MathModPi{\delta_{\tau_1}}{\delta_\chi}{\delta_{\tau_2}} }
    }
  }
}

\newcommand{\RepairTermOneInsPi}{
  \inferrule*
  [lab=\RInsPi]
  {
    {\turnstile
      {\dcontext{E}{\delta_E}{\Gamma}{\delta_\Gamma}}
      {\freshOne{\chi}{\MathSame}{\delta_\chi}}
    }
    \and
    {\MathPatches{\chi}{\delta_\chi}{x}}
    \and
    {\MathPatches{\tau}{\delta_{\tau_1}}{\tau_1}}
    \\\\\\
    \turnstile
    {\dcontext{E}{\delta_E}{\Gamma}
      {\MathIns{\MathLocalAssum{x}{\tau_1}}{\delta_\Gamma}}
    }
    { \repair{t}{\delta_{t}}{\tau}
      {\delta_{\tau_2}[\subst{\MathReplace{\chi}}{\MathReplace{x}}]}
    }
  }
  {
    \turnstile
    { \dcontext{E}{\delta_E}{\Gamma}{\delta_\Gamma} }
    { \repair
      {t}
      { \MathInsLam{x}{\delta_{t}} }
      {\tau}
      { \MathInsPi{\delta_{\tau_1}}{\chi}{\delta_{\tau_2}} }
    }
  }
}


\resizebox{\columnwidth}{!}{

$\turnstile
{ \dcontext{E}{\delta_E}{\Gamma}{\delta_\Gamma} }
{\repairTermWithoutType{t}{\delta_{t}}}$

% $\inferrule*
%   [lab=\UTRVar]
%   {
%     {\turnstile
%       { \dcontext{E}{\delta_E}{\Gamma}{\delta_\Gamma} }
%       { \declDiff{v}{\delta_v}{\delta_{\tau_v}} }
%     }
%     \and
%     {\turnstile
%       { \dcontext{E}{\delta_E}{\Gamma}{\delta_\Gamma} }
%       { \repairArgs{ [] }{\tau_v}{\delta_{\tau_v}}{\delta_v}{\delta}
%       }
%     }
%   }
%   {\turnstile
%     { \dcontext{E}{\delta_E}{\Gamma}{\delta_\Gamma} }
%     { \repairTermWithoutType{v}{\delta} }
%   }$

}

\end{document}

%%% Local Variables:
%%% mode: latex
%%% TeX-master: t
%%% End:

\section{Deriving repair functions}~\label{deriving}

Accounting for changes in inductive definitions requires quite a bit of work:

\begin{itemize}

\item the inductive type itself must be repaired, including its parameters, and
its indices,

\item each constructor must be repaired, including their parameters, and their
instantiation of indices,

\item the automatically-generated elimination principles for the type must also
be repaired accordingly.

\end{itemize}%
%
Let us focus on the latter.  In a proof assistant like Coq, when an inductive
data type \coqinline{t} is defined, the system also defines a family of
\textit{elimination principles} \coqinline{t_ind}, \coqinline{t_rec}, and
\coqinline{t_rect}, that essentially encode an induction principle for the given
type.  For instance, for the types \coqinline{list} and \coqinline{vec}, the
induction principles are given in Figure~\ref{list-ind-vec-ind}.

\begin{figure*}[!htp]

  \noindent%
  \begin{minipage}[t]{0.45\textwidth}
\begin{minted}[linenos=false,escapeinside=@@]{coq}
list_ind :
∀ (A : Type)          @\tikzmark{list1}@
  (P : list A →       @\tikzmark{list2}@
       Prop),
  P (nil A) →         @\tikzmark{list31}@
  (∀ (a : A)          @\tikzmark{list32}@
     (l : list A),
     P l →
     P (cons A a l)
  ) →
                      @\tikzmark{list4}@
  ∀ l : list A,       @\tikzmark{list5}@
    P l               @\tikzmark{list6}@
\end{minted}
\begin{tikzpicture}[remember picture,overlay]
  % \draw[help lines,xstep=.1,ystep=.1] (0,0) grid (8,10);
  % \foreach \x in {0,1,...,9} { \node [anchor=north] at (\x,0) {\x}; }
  % \foreach \y in {0,1,...,9} { \node [anchor=east] at (0,\y) {\y}; }
  \tikzstyle{circlednumber} = [shape=circle,draw,font=\bfseries,inner sep=1pt,thick];
  \node[circlednumber, above=-0.05 of pic cs:list1 ] {1};
  \node[circlednumber, above=-0.05 of pic cs:list2 ] {2};
  \node[circlednumber, above=-0.05 of pic cs:list31] {3};
  \node[circlednumber, above=-0.05 of pic cs:list32] {3};
  \node[circlednumber, above=-0.05 of pic cs:list4 ] {4};
  \node[circlednumber, above=-0.05 of pic cs:list5 ] {5};
  \node[circlednumber, above=-0.05 of pic cs:list6 ] {6};
\end{tikzpicture}
  \end{minipage}%
  \begin{minipage}[t]{0.55\textwidth}
\begin{minted}[linenos=false,escapeinside=@@]{coq}
vec_ind :
∀ (A : Type)                   @\tikzmark{vec1}@
  (P : ∀ n : nat, vec A n →    @\tikzmark{vec2}@
         Prop),
  P 0 (vnil A) →               @\tikzmark{vec31}@
  (∀ (a : A) (n : nat)         @\tikzmark{vec32}@
     (v : vec A n),
     P n v →
     P (S n) (vcons A a n v)
  ) →
  ∀ (n : nat)                  @\tikzmark{vec4}@
    (v : vec A n),             @\tikzmark{vec5}@
    P n v                      @\tikzmark{vec6}@
\end{minted}
\begin{tikzpicture}[remember picture,overlay]
  % \draw[help lines,xstep=.1,ystep=.1] (0,0) grid (8,10);
  % \foreach \x in {0,1,...,9} { \node [anchor=north] at (\x,0) {\x}; }
  % \foreach \y in {0,1,...,9} { \node [anchor=east] at (0,\y) {\y}; }
  \tikzstyle{circlednumber} = [shape=circle,draw,font=\bfseries,inner sep=1pt,thick];
  \node[circlednumber, above=-0.05 of pic cs:vec1 ] {1};
  \node[circlednumber, above=-0.05 of pic cs:vec2 ] {2};
  \node[circlednumber, above=-0.05 of pic cs:vec31] {3};
  \node[circlednumber, above=-0.05 of pic cs:vec32] {3};
  \node[circlednumber, above=-0.05 of pic cs:vec4 ] {4};
  \node[circlednumber, above=-0.05 of pic cs:vec5 ] {5};
  \node[circlednumber, above=-0.05 of pic cs:vec6 ] {6};
\end{tikzpicture}
  \end{minipage}

  \caption{ Induction principles for \coqinline{list} and \coqinline{vec}. }~\label{list-ind-vec-ind}

\end{figure*}

The property being returned \coqinline{P} is usually referred to as the
\textit{motive} of the elimination principle.  The motive is a property about a
value of the inductive type being eliminated: we will refer to this value as the
\textit{target} of the motive.  The value being eliminated (the last
universally quantified value, respectively \coqinline{l} and \coqinline{v}) are
usually referred to as the \textit{discriminee}.  The process that computes the
elimination principle's type from the inductive data type definition
$\Inductive{n}{\overrightarrow{p}}{\overrightarrow{i}}{u}{\overrightarrow{c}}$
is involved, but straightforward:

\begin{enumerate}[label=\protect\circled{\arabic*}]

\item~\label{ind-params} the inductive parameters $\overrightarrow{p}$ are
universally quantified over,

\item the motive \coqinline{P} is universally quantified over, its type being
computed by:

\begin{itemize}

\item universally quantifying over the inductive indices $\overrightarrow{i}$,

\item universally quantifying over a target, i.e.\ an element of the inductive
type undergoing elimination, fully-instantiated with parameters
from~\ref{ind-params} and the indices from the previous sub-step,

\item and returning in the appropriate universe for the desired elimination
principle,

\end{itemize}

\item for each constructor
$\Constructor{n_c}{\overrightarrow{p_c}}{\overrightarrow{i_c}}$, a case is
universally quantified over, whose type is obtained thus:

\begin{itemize}

\item each constructor parameter $\overrightarrow{p_c}$ is
universally quantified over, and for each parameter $p_{c_i}$ that is a
recursive occurrence of the data type being eliminated, a fully-instantiated
motive, with appropriate arguments so as to target this parameter, is also
universally quantified over immediately after this parameter,

\item the return type is also a fully-instantiated motive, where inductive
indices are instantiated with the constructor indices $\overrightarrow{i_c}$,
and the motive's target is an instantiation of the constructor $n_c$, with
inductive parameters from step~\ref{ind-params}, and constructor parameters from
the previous sub-step,

\end{itemize}

\item~\label{ind-indices} inductive indices $\overrightarrow{i}$ are
universally quantified over,

\item the discriminee is universally quantified over, its type being the
inductive type applied to inductive parameters from~\ref{ind-params} and
inductive indices from~\ref{ind-indices},

\item finally, the elimination rule returns a value of the fully-instantiated
motive, with indices from~\ref{ind-indices}, and the discriminee as target.

\end{enumerate}
%
We will use Haskell syntax to describe meta-language implementation details (we
reserve Coq syntax for programs in the object language).  Our function computing
the type of an eliminator has the following high-level skeleton:

\begin{minted}{haskell}
eliminatorType ... =
  quantifyVars  inductiveParameters                 {- 1 -}
. mkPi          motiveType            motive        {- 2 -}
. quantifyCases inductiveConstructors               {- 3 -}
. quantifyVars  inductiveIndices                    {- 4 -}
. mkPi          discrimineeType       discriminee   {- 5 -}
$ mkApp         indices               discriminee   {- 6 -}
  where
   indices = applyVars inductiveIndices motive
\end{minted}

\noindent where \coqinline{quantifyVariables}, \coqinline{applyVariables} are
folding functions that create \coqinline{Pi}s, and \coqinline{App}s,
respectively, and \coqinline{motiveType} and \coqinline{discrimineeType} are
computed by additional \coqinline{fold} operations.

Our insight is as follows: while we cannot derive a diff-propagating function
for any \coqinline{fold} in general, we can describe how a given folding
operation \coqinline{f} reacts to differences in the input list \coqinline{l} in
the call \coqinline{(fold f l b)}.  From this description, we can derive the
diff-propagating function for the whole \coqinline{fold} operation.  To make
things more concrete, let's focus on \coqinline{quan\-tify\-Variables} and build
its corresponding, diff-propagating version \coqinline{δquan\-tify\-Variables}.
The original function is a right-fold that universally quantifies over every
element encountered. We can describe how it reacts to changes in the input list
using the following record type:

\begin{minted}{haskell}
data ΔListFold τ δτ δ = ΔListFold
  { onInsert  ::     τ →     [τ] → δ → δ
  , onModify  ::    δτ → τ → [τ] → δ → δ
  , onPermute :: [Int] →     [τ] → δ → δ
  , onRemove  ::         τ → [τ] → δ → δ
  , onReplace ::   [τ] →     [τ] → δ → δ
  , onSame    ::             [τ] → δ → δ
  }
\end{minted}

\noindent For each diff operation on lists, we record a transformer on the
output diff type \coqinline{δ}.  In our example, the fold is building a
\coqinline{Term}, so the output diff type is \coqinline{ΔTerm}, and we describe
how the output term is affected by changes to the input list: when an element is
inserted in the input list, a $\Pi$ is added to the output term, when an element
is removed from the input list, a $\Pi$ is removed from the output term, etc.
When the entire list gets replaced, the handler receives the old list
\coqinline{l} and the new list \coqinline{l'}: the output term will see
\coqinline{(length l')} $\Pi$s removed, and for each element of \coqinline{l}, a
corresponding $\Pi$ will be added.  We then have two functions:

\begin{minted}{haskell}
δListFoldLeft, δListFoldRight ::
  ΔListFold τ δτ δ ->
  [τ] -> ΔList τ δτ -> Maybe (δ -> δ)
\end{minted}

\noindent which take such a description, an original list, and a list diff, and
compute the resulting diff transformer for the output type, as long as the diff
makes sense for the list.  Receiving the original list lets us check the sanity
of the list diff at each step (e.g.\ a diff should not ask to drop the head from
an empty list), and lets us pass the element being modified or removed to the
handlers who need the information.  With all this machinery ready, we can derive
\coqinline{δEliminator\-Type} as:

\begin{minted}{haskell}
δEliminatorType ... =
    δquantifyVars  ips              δips
<$> CopyPi         δmotiveType      Same
<$> δquantifyCases cs               δcs
 $  δquantifyVars  iis              δiis
 $  CopyPi         δdiscrimineeType Same
 $  CopyApp        δindices         Same
  where
    δindices = δapplyVars iis δiis Same
\end{minted}

\noindent which matches closely the original definition of
\coqinline{Eliminator\-Type}.  We use the same mechanism to derive a
diff-propagating version of the function that computes the type of an inductive
type family (based on the diff of its parameters and indices lists), and a
diff-propagating version of the function that computes the type of a constructor
(based on the diff of its parameters and indices lists).  At a high-level, we
have the following property:

\[
  {
    \inferrule*
    [width=5em]
    {
      {
        {\MathMkElimType
          {\Inductive
            {n}
            {\overrightarrow{p}}
            {\overrightarrow{i}}
            {u}
            {\overrightarrow{c}}
          }
        }
        \op{\ =\ }
        \out{e}
      }
      \\
      {\MathPatches
        {\Inductive
          {n}
          {\overrightarrow{p}}
          {\overrightarrow{i}}
          {u}
          {\overrightarrow{c}}
        }
        {\delta_i}
        {\Inductive
          {n'}
          {\overrightarrow{p'}}
          {\overrightarrow{i'}}
          {u'}
          {\overrightarrow{c'}}
        }
      }
      \\
      {
        {\MathMkΔElimType
          {\Inductive
            {n}
            {\overrightarrow{p}}
            {\overrightarrow{i}}
            {u}
            {\overrightarrow{c}}
          }
          {\delta_i}
        }
        \op{\ =\ }
        \out{\delta_e}
      }
      \\
      {
        {\MathMkElimType
          {\Inductive
            {n'}
            {\overrightarrow{p'}}
            {\overrightarrow{i'}}
            {u'}
            {\overrightarrow{c'}}
          }
        }
        \op{\ =\ }
        \out{e'}
      }
    }
    {
      \MathPatches{e}{\delta_e}{e'}
    }
  }
\]

\noindent i.e., \coqinline{δEliminatorType} produces the correct diff to
transport the old eliminator type to the new eliminator type.  In fact, we even
have the following corollary:

\[
  {
    \inferrule*
    [width=5em]
    {
      {{\MathMkElimType{i}} \op{\ =\ } \out{e}}
      \\
      {
        {\MathMkΔElimType
          {i}
          {\ModifyInductive
            {\MathReplace{n}}
            {\MathReplace{\overrightarrow{p}}}
            {\MathReplace{\overrightarrow{i}}}
            {\MathReplace{u}}
            {\MathReplace{\overrightarrow{c}}}
          }
        }
        \op{\ =\ }
        \out{\delta_e}
      }
      \\
      {
        {\MathMkElimType
          {\Inductive
            {n}
            {\overrightarrow{p}}
            {\overrightarrow{i}}
            {u}
            {\overrightarrow{c}}
          }
        }
        \op{\ =\ }
        \out{e'}
      }
    }
    {\MathPatches{e}{\delta_e}{e'}}
  }
\]

\noindent making \coqinline{eliminatorType} redundant as it coincides with a
modification where each field of the inductive data type definition is replaced
with a new value.

\documentclass[preview]{standalone}

\usepackage{fontspec}
\usepackage{gfsdidot}          % makes math stuff look neat, but also overrides text stuff
\setmainfont{TeX Gyre Pagella} % forces back text stuff to look neat too
\setmonofont{DejaVuSansMono}   % monospace stuff should look neat too

\usepackage[dvipsnames]{xcolor} % must be loaded prior to tikz, pgfplots, etc.
\let\iint\relax
\let\iiint\relax
\let\iiiint\relax
\let\idotsint\relax
\usepackage{amsmath} % loaded by mathtools?
\usepackage{amssymb}
\usepackage[english]{babel}
\usepackage{booktabs}
\usepackage{calc}
\usepackage{caption}
\usepackage{color}
\usepackage{DejaVuSansMono}
\usepackage{dsfont}
\let\textlozenge\relax
\usepackage[inline]{enumitem}
\usepackage{environ}
\usepackage{etoolbox}
\usepackage{fontawesome}
\usepackage{forest}
\usepackage{graphicx}
\usepackage{longtable}
\usepackage{ltablex} \keepXColumns % Made me sad for appendix tables, where was it used?
\usepackage{makecell}
\usepackage{mathpartir}
%\usepackage{mathtools} % \coloneqq
\usepackage{mdframed}
\usepackage{minted}
\usepackage{multicol}
\usepackage[parfill]{parskip}
\usepackage{pgfplots}
\usepackage{pifont} % \ding
\usepackage{soul}
\usepackage{syntax}
\usepackage{tabularx}
\usepackage[skins,theorems]{tcolorbox}
\usepackage{tikz}
\usepackage{tikzpeople}
\usepackage{tikz-qtree}
\usepackage{xpatch}
\usepackage[breaklinks=true,pdfborder={0 0 0}]{hyperref}

% gfsdidot makes ∀ and ∃ super ugly, this reverts them to something fine
\DeclareSymbolFont{CMsymbols}{OMS}{cmsy}{b}{n}
\SetSymbolFont{CMsymbols}{bold}{OMS}{cmsy}{ub}{n}
\DeclareMathSymbol{\forall}{\mathord}{CMsymbols}{"38}
\DeclareMathSymbol{\exists}{\mathord}{CMsymbols}{"39}

\definecolor{monokaibg}{HTML}{272822}
\definecolor{color01}{HTML}{E6194B}
\definecolor{color02}{HTML}{3CB44B}
\definecolor{color03}{HTML}{FFE119}
\definecolor{color04}{HTML}{4363D8}
\definecolor{color05}{HTML}{F58231}
\definecolor{color06}{HTML}{911EB4}
\definecolor{color07}{HTML}{46F0F0}
\definecolor{color08}{HTML}{F032E6}
\definecolor{color09}{HTML}{BCF60C}
\definecolor{color10}{HTML}{FABEBE}
\definecolor{color11}{HTML}{008080}
\definecolor{color12}{HTML}{E6BEFF}
\definecolor{color13}{HTML}{9A6324}
\definecolor{color14}{HTML}{FFFAC8}
\definecolor{color15}{HTML}{800000}
\definecolor{color16}{HTML}{AAFFC3}
\definecolor{color17}{HTML}{808000}
\definecolor{color18}{HTML}{FFD8B1}
\definecolor{color19}{HTML}{000075}
\definecolor{color20}{HTML}{808080}
\definecolor{color21}{HTML}{FFFFFF}
\definecolor{color22}{HTML}{000000}

\def\mathunderline#1#2{\color{#1}\underline{{\color{black}#2}}\color{black}}

\forestset{
  default preamble={
    for tree={
      align=center,
      draw=black,
      edge={
        ultra thick,
      },
      edge path={
        \noexpand\path[\forestoption{edge}]
        (!u.parent anchor)
        -- +(0,-10pt)
        -| (.child anchor)
        \forestoption{edge label};
      },
      font=\bfseries,
      % inner sep=10pt,
      l sep=30pt,
      % line width=2pt,
      minimum height=20pt,
      minimum width=30pt,
      parent anchor=south,
      rectangle,
      ultra thick,
    }
  }
}

\newcolumntype{Y}{>{\centering\arraybackslash}X}

% defines safecoqinline, a command like coqinline but works in tabular
% environments
\makeatletter
\def\safecoqinline#1{%
\ifx\@footnotetext\TX@trial@ftn
\detokenize{#1}%
\else
\coqinline{#1}%
\fi}
\makeatother

\captionsetup{
  justification=centering,
}

\usetikzlibrary{arrows, calc, fit, positioning, shapes, tikzmark}

\tikzset{
bicolor/.style 2 args={
  dashed,dash pattern=on 4pt off 4pt,#1,
  postaction={draw,dashed,dash pattern=on 4pt off 4pt,#2,dash phase=4pt}
  },
}

\tikzstyle{Matching}=[
black,
dashed,
line width=3pt,
]

\tikzstyle{RoundedDottedPath}=[
densely dotted,
color05,
line width=3pt,
rounded corners=10pt,
]

\tikzstyle{RoundedRectangle}=[
dashed,
draw,
line width=3pt,
rounded corners=10pt,
]

\tikzstyle{NodeLabel}=[
black,
circle,
draw,
fill=white,
font=\bfseries,
inner sep=1pt,
ultra thick,
]

\graphicspath{ {./images/} }

% Fix a mdframed bug where skipbelow is ignored
% \makeatletter
% \xpatchcmd{\endmdframed}
%   {\aftergroup\endmdf@trivlist\color@endgroup}
%   {\endmdf@trivlist\color@endgroup\@doendpe}
%   {}{}
% \makeatother

% \surroundwithmdframed[
% backgroundcolor=monokaibg,
% linecolor=white,
% skipabove=1em,
% skipbelow=1em,
% leftmargin=10pt,
% innertopmargin=1pt,
% innerbottommargin=0pt
% ]{minted}

\setminted{
  % DO NOT use bgcolor, use mdframed instead
  % bgcolor=monokaibg,
  linenos=true,
  style=lovelace,
  breaklines=true,
  encoding=utf8,
  fontsize=\large,
  baselinestretch=1,
}

% need bgcolor for the display to be nice in rules
\newmintinline{coq}{bgcolor=white}

\makeatletter
\pgfdeclareshape{document}{
\inheritsavedanchors[from=rectangle] % this is nearly a rectangle
\inheritanchorborder[from=rectangle]
\inheritanchor[from=rectangle]{center}
\inheritanchor[from=rectangle]{north}
\inheritanchor[from=rectangle]{south}
\inheritanchor[from=rectangle]{west}
\inheritanchor[from=rectangle]{east}
% ... and possibly more
\backgroundpath{% this is new
% store lower right in xa/ya and upper right in xb/yb
\southwest \pgf@xa=\pgf@x \pgf@ya=\pgf@y
\northeast \pgf@xb=\pgf@x \pgf@yb=\pgf@y
% compute corner of ‘‘flipped page’’
\pgf@xc=\pgf@xb \advance\pgf@xc by-10pt % this should be a parameter
\pgf@yc=\pgf@yb \advance\pgf@yc by-10pt
% construct main path
\pgfpathmoveto{\pgfpoint{\pgf@xa}{\pgf@ya}}
\pgfpathlineto{\pgfpoint{\pgf@xa}{\pgf@yb}}
\pgfpathlineto{\pgfpoint{\pgf@xc}{\pgf@yb}}
\pgfpathlineto{\pgfpoint{\pgf@xb}{\pgf@yc}}
\pgfpathlineto{\pgfpoint{\pgf@xb}{\pgf@ya}}
\pgfpathclose
% add little corner
\pgfpathmoveto{\pgfpoint{\pgf@xc}{\pgf@yb}}
\pgfpathlineto{\pgfpoint{\pgf@xc}{\pgf@yc}}
\pgfpathlineto{\pgfpoint{\pgf@xb}{\pgf@yc}}
\pgfpathlineto{\pgfpoint{\pgf@xc}{\pgf@yc}}
}
}
\makeatother

%%%%% General-purpose text macros %%%%%
\newcommand{\define}[1]{\emph{#1}}
\newcommand{\mycite}[1]{\citeauthor{#1}~\cite{#1}}

\newcommand{\Language}[1]{\emph{#1}}
\newcommand{\Tool}[1]{\emph{#1}}
\newcommand{\Gallina}{\Language{Gallina}}
\newcommand{\Ltac}{\Language{Ltac}}
\newcommand{\Vernacular}{\Language{Vernacular}}

\newcommand{\Chick}{\Language{Chick}}
\newcommand{\Coq}{\Language{Coq}}
\newcommand{\CoqIDE}{\Tool{CoqIDE}}
\newcommand{\Haskell}{\Language{Haskell}}
\newcommand{\JavaScript}{\Language{JavaScript}}
\newcommand{\OCaml}{\Language{OCaml}}
\newcommand{\PeaCoq}{\Language{PeaCoq}}
\newcommand{\RxJS}{\Language{RxJS}}
\newcommand{\SerAPI}{\Tool{SerAPI}}
\newcommand{\Snap}{\Language{Snap}}
\newcommand{\TypeScript}{\Language{TypeScript}}

\newcommand{\OperatorColor}{purple}
\newcommand{\Operator}[1]{\textcolor{\OperatorColor}{\ #1\ }}
\newcommand{\Entails}{\Operator{\vdash}}
\newcommand{\HasType}{\Operator{:}}

%\newmintinline{coq}{fontsize=\small}
% \newcommand{\coqinline}[1]{%
%   %\colorbox{monokaibg}{%
%   \parbox[c][0.9em]{\widthof{\mycoq{#1}}}{\mycoq{#1}}%
%   %}%
% }

\newcommand{\modified}[1]{\color{Orange}{#1}}
\newcommand{\repaired}[1]{\color{PineGreen}{#1}}

\newcommand{\rulename}[1]{$\LeftTirNameStyle{#1}$}

\newcommand{\RmApp}{Rm-App}
\newcommand{\RmPi}{Rm-Pi}
\newcommand{\InsApp}{Ins-App}

%%%%% Math-mode macros %%%%%

\newcommand{\opcolor}{purple}
\newcommand{\out}[1]{ \boxed{ \textcolor{teal}{#1} } }
\newcommand{\MathPatches}[3]{%
#1 \overset{#2}{\mathbin{\textcolor{\opcolor}{\rightsquigarrow}}} \out{#3}%
}

\newcommand{\App}{\$}
\newcommand{\Mod}{Mod}
\newcommand{\Drop}{Drop}
\newcommand{\Ins}{Ins}
\newcommand{\Keep}{Keep}
\newcommand{\Lam}{\lambda}

\newcommand{\oMod}{\overset{\mathtt{\Mod}}}
\newcommand{\oDrop}{\overset{\mathtt{\Drop}}}
\newcommand{\oIns}{\overset{\mathtt{\Ins}}}
\newcommand{\oKeep}{\overset{\mathtt{\Keep}}}

\newcommand{\Cons}{::}
\newcommand{\permute}[2]{%
\overline{#2}^{\overset{#1}{\rightleftarrows}}%
}
\newcommand{\permuteOp}[2]{%
\overset{\overset{#1}{\rightleftarrows}}{#2}%
}

\newcommand{\mkMathPiRaw}[4]{#3{\Pi} #4{#1} \rightarrow #2}
\newcommand{\mkMathPi}[5]{\mkMathPiRaw{(#2 : #1)}{#3}{#4}{#5}}


\newcommand{\MathCons}[2]{#1 \Cons #2}
\newcommand{\MathDrop}[1]{\oDrop{\Cons} #1}
\newcommand{\MathDropApp}[1]{\oDrop{\App} #1}
\newcommand{\MathDropLam}[1]{\oDrop{\Lam} #1}
\newcommand{\MathDropPi}[1]{\oDrop{\Pi} #1}
\newcommand{\MathHole}{\texttt{\_}}
\newcommand{\MathIns}[2]{#1 \oIns{\Cons} #2}
\newcommand{\MathInsApp}[2]{\mkMathApp{#1}{#2}{\oIns}}
\newcommand{\MathInsAppOp}{\oIns{\App}}
\newcommand{\MathInsLam}[2]{\mkMathLam{#1}{#2}{\oIns}{}}
\newcommand{\MathInsLamOp}{\oIns{\Lam}}
\newcommand{\MathInsPi}[3]{\mkMathPi {#1}{#2}{#3}{\oIns}{}}
\newcommand{\MathInsPiOp}{\oIns{\Pi}}
\newcommand{\MathLam}[2]{\mkMathLam{#1}{#2}{}{}}
\newcommand{\MathLams}[3]{\mkMathLam{#1}{#2}{}{\BarCount{#3}}}
\newcommand{\MathKeepPiOp}{\oKeep{\Pi}}
\newcommand{\MathMod}[2]{#1 \oMod{\Cons} #2}
\newcommand{\MathModAppOp}{\oMod{\App}}
\newcommand{\MathModLamOp}{\oMod{\Lam}}
\newcommand{\MathModPiOp}{\oMod{\Pi}}
\newcommand{\MathPermute}[2]{\permuteOp{#1}{\Cons} #2}
\newcommand{\MathPermuteLams}[2]{\permuteOp{#1}{\lambda} #2}
\newcommand{\MathPermutePis}[2]{\permuteOp{#1}{\Pi} #2}
\newcommand{\MathPi}[3]{\mkMathPi{#1}{#2}{#3}{}{}}
\newcommand{\MathPis}[4]{\mkMathPi{#1}{#2}{#3}{}{\BarCount{#4}}}
\newcommand{\MathReplace}[1]{\mathds{K}({#1})}
\newcommand{\MathSame}{\mathds{1}}

\newcommand{\mkMathApp}[3]{#1 #3{ \$ } #2}

%%% REPAIR %%%

\newcommand{\op}[1]{\textcolor{\opcolor}{#1}}

\newcommand{\boundOp}{\op{\text{Bound}}}
\newcommand{\freshOneOp}{\op{\text{Fresh}_1}}
\newcommand{\freshTwoOp}{\op{\text{Fresh}_2}}
\newcommand{\repairIndOp}{\op{R_{I}}}
\newcommand{\repairProgOp}{\op{R_{P}}}
\newcommand{\repairTermOneOp}{\op{R_{T_1}}}
\newcommand{\repairTermTwoOp}{\op{R_{T_2}}}
\newcommand{\repairTermThreeOp}{\op{R_{T_2}}}
\newcommand{\repairVernacOp}{\op{R_{V}}}
\newcommand{\vernacEnvOp}{\op{E_{V}}}
\newcommand{\repairBranchesOp}{\op{R_{B}}}

\newcommand{\blackbrackets}[1]{\left[ #1 \right]}
\newcommand{\blackdiff}[2]{\blackbrackets{\genfrac{}{}{0pt}{}{#1}{#2}}}
\newcommand{\bound}[1]{\boundOp\op{(}#1\op{)}}
\newcommand{\brackets}[1]{%
  \color{\opcolor} \left[ \normalcolor #1 \color{\opcolor} \right] \normalcolor%
}
\newcommand{\context}[2]{#1 \op{,} #2}
\newcommand{\dcontext}[4]{\context{\diff{#1}{#2}}{\diff{#3}{#4}}}
\newcommand{\denv}[2]{\diff{#1}{#2}}
\newcommand{\diff}[2]{\brackets{\genfrac{}{}{0pt}{}{#1}{#2}}}
\newcommand{\dtau}[1]{\delta_{\tau_{#1}}}
\newcommand{\equals}[2]{#1 \mathbin{\textcolor{\opcolor}{=}} #2}
\newcommand{\freshOne}[3]{\freshOneOp\op{(}\diff{#1}{#2}\op{) =}\ \out{#3}}
\newcommand{\freshTwo}[3]{\freshTwoOp\op{(}#1\op{,}#2\op{) =}\ \out{#3}}
\newcommand{\genericrepair}[3]{\mkrepair{\repairTermTwoOp}{#1}{#2}{#3}{\op{?}}}
\newcommand{\nturnstile}[2]{#1 \mathbin{\textcolor{\opcolor}{\nvdash}} #2}
\newcommand{\qmark}{\textcolor{\opcolor}{?}}
\newcommand{\squiggly}[2]{#1 \mathbin{\textcolor{\opcolor}{\rightsquigarrow}} #2}
\newcommand{\subst}[2]{#1 \leftarrow #2}
\newcommand{\turnstile}[2]{#1\ \mathbin{\textcolor{\opcolor}{\vdash}}\ #2}
\newcommand{\repair}[4]{ \mkrepair{\repairTermOneOp}{#1}{#2}{#3}{#4} }
\newcommand{\repairBranches}[3]{\mksimplerepair{\repairBranchesOp}{#1\op{,}\ #2}{#3}}
\newcommand{\repairInd}[7]{
  \mksimplerepair%
  {\repairIndOp}
  {\diff{#1}{#2}\op{,}\ \diff{#3}{#4}\op{,}\ \diff{#5}{#6}}
  {#7}
}
\newcommand{\repairProg}[3]{
  \repairProgOp\op{(}\blackdiff{#1}{#2}\op{)}\ \op{=}\ \out{#3}
}
\newcommand{\repairTerm}[4]{\mkrepair{\repairTermOneOp}{#1}{#2}{#3}{#4}}
\newcommand{\repairTermWithoutType}[2]{\mksimplerepair{\repairTermThreeOp}{#1}{#2}}
\newcommand{\repairVernac}[3]{
  \repairVernacOp\op{(}\blackdiff{#1}{#2}\op{)}\ \op{=}\ \out{#3}
}

\newcommand{\DefinitionText}{\text{Definition}}
\newcommand{\InductiveText}{\text{Inductive}}
\newcommand{\ConstructorText}{\text{Constructor}}
\newcommand{\GlobalDefinitionText}{\text{GlobalDefinition}}
\newcommand{\GlobalInductiveText}{\text{GlobalInductive}}

\newcommand{\Definition}[4]{ \DefinitionText(\{ #1, #2, #3, #4 \}) }
\newcommand{\Inductive}[5]{ \InductiveText(\{ #1, #2, #3, #4, #5 \}) }
\newcommand{\Constructor}[3]{ \ConstructorText(\{ #1, #2, #3 \}) }
\newcommand{\GlobalDefinition}[3]{ \GlobalDefinitionText(\{ #1, #2, #3 \}) }
\newcommand{\GlobalInductive}[5]{ \GlobalInductiveText(\{ #1, #2, #3, #4, #5 \}) }

\newcommand{\ModifyDefinition}[4]{ \delta_{\DefinitionText}(\{ #1, #2, #3, #4 \}) }
\newcommand{\ModifyInductive}[5]{ \delta_{\InductiveText}(\{ #1, #2, #3, #4, #5 \}) }
\newcommand{\ModifyGlobalDefinition}[3]{ \delta_{\GlobalDefinitionText}(\{ #1, #2, #3 \}) }
\newcommand{\ModifyGlobalInductive}[5]{ \delta_{\GlobalInductiveText}(\{ #1, #2, #3, #4, #5 \}) }

\newcommand{\GlobalDefinitionAnon}{\GlobalDefinitionText(\{ \ldots \})}
\newcommand{\GlobalInductiveAnon}{\GlobalInductiveText(\{ \ldots \})}

\newcommand{\DefinitionAnon}{\DefinitionText(\{ \ldots \})}
\newcommand{\InductiveAnon}{\InductiveText(\{ \ldots \})}

\newcommand{\RepairProg}[1]{R-Prog-#1}

\newcommand{\RProgMod}{\RepairProg{Modify}}
\newcommand{\RProgKeep}{\RepairProg{Keep}}
\newcommand{\RProgSameCons}{\RepairProg{Same-Cons}}

\newcommand{\scalefactor}{1}
\newcommand{\invscalefactor}{1}

\NewEnviron
    {Rules}[2]
    {
      \begin{figure*}
        %\scalebox{\invscalefactor}{
          %\begin{minipage}{\scalefactor\textwidth}
            %\begin{mathpar}
              \BODY
            %\end{mathpar}
          %\end{minipage}
        %}
        \caption{#2}
        \label{#1}
      \end{figure*}
    }

\newcommand{\MathSameProg}{\MathSame_P}
\newcommand{\MathSameVernac}{\MathSame_V}
\newcommand{\MathSameGlobal}{\MathSame_G}

\newcommand{\MathProp}{\mathtt{Prop}}
\newcommand{\MathSet} {\mathtt{Set}}
\newcommand{\MathType}{\mathtt{Type}}

\newcommand{\hasType}[2]{ #1\ \op{:}\ #2 }

\newcommand{\mkrepair}[5]{
  #1
  \op{(}
  \hasType
  { #2 }
  { \diff{ #4 }{ #5 } }
  \op{) =\ }
  \out{ #3 }
}

\newcommand{\MathLocalAssum}[2]{ ( #1 : #2 ) }
\newcommand{\MathLocalDef}[3]{ ( #1 : #2 \coloneqq #3 ) }

\newcommand{\mksimplerepair}[3]{
  #1
  \op{(}
  #2
  \op{) =\ }
  \out{ #3 }
}

\newcommand{\RModPi}{\RepairTerm{Mod-$\Pi$}}
\newcommand{\RInsPi}{\RepairTerm{Ins-$\Pi$}}
\newcommand{\RDropPi}{\RepairTerm{Drop-$\Pi$}}
\newcommand{\RPermutePis}{\RepairTerm{Permute-$\Pi$s}}
\newcommand{\RReplace}{\RepairTerm{Replace}}
\newcommand{\RSame}{\RepairTerm{$\MathSame{}$-Other}}
\newcommand{\RSamePi}{\RepairTerm{$\MathSame{}$-$\Pi$}}

\newcommand{\RepairTermPrefix}{RT1}
\newcommand{\GenericRepairPrefix}{R-Term-2}
\newcommand{\UnknownTypeRepairPrefix}{R-Term-2}

\newcommand{\RepairTerm}[1]{\RepairTermPrefix-#1}
\newcommand{\GenericRepair}[1]{\GenericRepairPrefix-#1}
\newcommand{\UnknownTypeRepair}[1]{\UnknownTypeRepairPrefix-#1}

\newcommand{\mkMathLam}[4]{#3{\lambda} #4{#1} \rightarrow #2}

\newcommand{\MathModApp}[2]{ \mkMathApp{#1}{#2}{\oMod}{} }
\newcommand{\MathModLam}[2]{ \mkMathLam{#1}{#2}{\oMod}{} }
\newcommand{\MathModPi} [3]{ \mkMathPi{#1}{#2}{#3}{\oMod}{} }
\newcommand{\MathKeepPi}[1]{ \MathKeepPiOp \rightarrow #1 }

\newcommand{\MathAnnot}[2]{#1 : #2}

\newcommand{\UTRApp}{\UnknownTypeRepair{App}}
\newcommand{\UTRMatch}{\UnknownTypeRepair{Match}}
\newcommand{\UTRPi}{\UnknownTypeRepair{Pi}}
\newcommand{\UTRType}{\UnknownTypeRepair{Universe}}
\newcommand{\UTRVar}{\UnknownTypeRepair{Var}}
\newcommand{\UTRAnnot}{\UnknownTypeRepair{Annot}}
\newcommand{\UTRHole}{\UnknownTypeRepair{Hole}}
\newcommand{\UTROtherwise}{\UnknownTypeRepair{Otherwise}}

\newcommand{\repairArgsOp}{\op{R_{Fn}}}

\newcommand{\repairArgs}[5]{
  \op{\repairArgsOp(}#1\op{,} #2\op{,} #3\op{,} #4\op{) =}\ \out{#5}
}

\newcommand{\declDiff}[3]{
  \hasType
      { \diff{ #1 }{ \out{#2} } }
      { \diff{ \op{?} }{ \out{#3} } }
}

\newcommand{\MathMatch}[2]{ \text{match}\ #1\ \text{with}\ #2 }


\begin{document}
\begin{figure}[htp!]
  \resizebox{\columnwidth}{!}{

\begin{forest}
[,phantom
  [→,fill=color05,
    [→,fill=color04,
      tikz={
        \filldraw[RoundedDottedPath,color04,fill opacity=0.5]%,draw opacity=0.75]
        (.south west) -- (!l.south west)
        -- (!l.south east) -- (!l.north east)
        -- (.north east) -- (.north west) -- (.south west)
        ;
        % \draw[RoundedDottedPath,bicolor={color01}{color01}]
        % (.south east) -- (.north east) -- (.north west) -- (!1.north west) -- (!1.south west)
        % -- (!1.south east) -- (.south east)
        % ;
      },
      [A,fill=color01]
      [→,
      fill=color04,
      % tikz={\node[RoundedRectangle,color04,fit= () (!u)]{};},
      tikz={
          % \draw[RoundedDottedPath,bicolor={color02}{color02}]
          % (.south east) -- (.north east) -- (.north west) -- (!1.north west) -- (!1.south west)
          % -- (!1.south east) -- (.south east)
          % ;
        },
        [C,fill=color02]
        [D,fill=color03]
      ]
    ]
    [→,fill=color05,
      [A,fill=color12]
      [→,fill=color05,
      % tikz={\node[RoundedRectangle,color05,fit=()(!u)(!uu)]{};},
      tikz={\filldraw[RoundedDottedPath,color05,fill opacity=0.5]
          (.north east) -- (!u.north east) -- (!uu.north east)
          -- (!uu.north west) -- (!uu.south west) -- (!u.south west)
          -- (.south west) -- (.south east) -- (.north east)
          ;},
        [C,fill=color13]
        [D,fill=color14]
      ]
    ]
  ]
  [,phantom[,phantom[,phantom]]]
  [→,fill=color05,
    [→,fill=color04,
        tikz={
          \filldraw[RoundedDottedPath,color04,fill opacity=0.5]%,draw opacity=0.75]
          (.south west) -- (!l.south west) -- (!ll.south west)
          -- (!ll.south east) -- (!ll.north east) -- (!l.north east)
          -- (.north east) -- (.north west) -- (.south west)
          ;
          % \draw[RoundedDottedPath,bicolor={color02}{color02}]
          % (.south east) -- (.north east) -- (.north west) -- (!1.north west) -- (!1.south west)
          % -- (!1.south east) -- (.south east)
          % ;
        },
      [C,fill=color02]
      [→,fill=color04,
        [B]
        [→,fill=color04,
        tikz={
          % \draw[RoundedDottedPath,bicolor={color01}{color01}]
          % (.south east) -- (.north east) -- (.north west) -- (!1.north west) -- (!1.south west)
          % -- (!1.south east) -- (.south east)
          % ;
          },
        % tikz={\node[RoundedRectangle,color04,fit=()(!u)(!uu)]{};},
        % tikz={
        %   \draw[RoundedDottedPath,color04]
        %   (.north east) -- (!u.north east) -- (!uu.north east)
        %   -- (!uu.north west) -- (!uu.south west) -- (!u.south west)
        %   -- (.south west) -- (.south east) -- (.north east)
        %   ;
        % },
          [A,fill=color01]
          [D,fill=color03]
        ]
      ]
    ]
    [→,fill=color05,
      [A,fill=color12]
      [→,fill=color05,
        [B]
        [→,fill=color05,
        % tikz={\node[RoundedRectangle,color05,fit=()(!u)(!uu)(!uuu)]{};},
        tikz={\filldraw[RoundedDottedPath,color05,fill opacity=0.5]
          (.north east) -- (!u.north east) -- (!uu.north east) -- (!uuu.north east)
          -- (!uuu.north west) -- (!uuu.south west) -- (!uu.south west) -- (!u.south west)
          -- (.south west) -- (.south east) -- (.north east)
          ;},
          [C,fill=color13]
          [D,fill=color14]
        ]
      ]
    ]
  ]
]
\end{forest}


  }
\end{figure}
\end{document}

%%% Local Variables:
%%% mode: latex
%%% TeX-master: t
%%% End:

